\begin{quadro}[ht!]
  \centering\scriptsize
\caption{Unidade Curricular Algoritmos 1}

\begin{tabular}{|p{2.5cm}|p{2.5cm}|p{2.5cm}|p{2.5cm}|p{2.5cm}|}\hline

  Unidade curricular & \multicolumn{4}{p{10cm}|}{ Algoritmos 1 }\\\hline
  Código & 1.2 & Período & \multicolumn{2}{p{5cm}|}{ 1 }\\\hline
  Pré-Requisitos & \multicolumn{4}{p{10cm}|}{  } \\\hline
  \cellcolor{blue1}{ CH PRESENCIAL } &
    \cellcolor{blue1}{CH EAD Síncrona} &
    \cellcolor{blue1}{CH EAD Síncrona Mediada} &
    \cellcolor{blue1}{CH EAD Assíncrona} &
    \cellcolor{blue1}{CH APCC} \\\hline
  \cellcolor{blue1}{ 60 } &
    \cellcolor{blue1}{ 0 } &
    \cellcolor{blue1}{ 0 } &
    \cellcolor{blue1}{ 30 } &
    \cellcolor{blue1}{0} \\\hline
  \cellcolor{blue1}{ CH Teórica } &
    \cellcolor{blue1}{CH Prática} &
    \cellcolor{blue1}{CH Humanidades} &
    \cellcolor{blue1}{CH Extensão} &
    \cellcolor{blue1}{CH EAD TOTAL} \\\hline
  \cellcolor{blue1}{ 45 } &
    \cellcolor{blue1}{ 45 } &
    \cellcolor{blue1}{ 0 } &
    \cellcolor{blue1}{ 0 } &
    \cellcolor{blue1}{ 30 } \\\hline
  \cellcolor{blue1}{CARGA-HORÁRIA TOTAL} & \multicolumn{4}{p{10cm}|}{ \cellcolor{blue1}{ 90 } } \\\hline
  Área de conhecimento & \multicolumn{4}{p{10cm}|}{ Computação } \\\hline
  Ementa & \multicolumn{4}{p{10cm}|}{ Computadores: organização, características, funcionamento, descrição narrativa e fluxogramas. Algoritmos: tipos primitivos de dados, constantes e variáveis, expressões lógicas e aritméticas e práticas com programação. Estrutura de controle: desvio condicional e laços de repetição e práticas com programação. Modularização: funções, procedimentos, passagem de parâmetros e práticas com programação. Vetores, matrizes e registros: definições, manipulação e práticas com programação. Fundamentos de teste de software: teste de mesa e depuração. } \\\hline
  Competências & \multicolumn{4}{p{10cm}|}{ C1 } \\\hline


  
	\end{tabular}

\end{quadro}


\clearpage
\newpage

\begin{quadro}[ht!]
  \centering\scriptsize
\caption{Unidade Curricular Interpretação e Produção de Textos Científicos e Profissionais}

\begin{tabular}{|p{2.5cm}|p{2.5cm}|p{2.5cm}|p{2.5cm}|p{2.5cm}|}\hline

  Unidade curricular & \multicolumn{4}{p{10cm}|}{ Interpretação e Produção de Textos Científicos e Profissionais }\\\hline
  Código & 1.3 & Período & \multicolumn{2}{p{5cm}|}{ 1 }\\\hline
  Pré-Requisitos & \multicolumn{4}{p{10cm}|}{  } \\\hline
  \cellcolor{blue1}{ CH PRESENCIAL } &
    \cellcolor{blue1}{CH EAD Síncrona} &
    \cellcolor{blue1}{CH EAD Síncrona Mediada} &
    \cellcolor{blue1}{CH EAD Assíncrona} &
    \cellcolor{blue1}{CH APCC} \\\hline
  \cellcolor{blue1}{ 30 } &
    \cellcolor{blue1}{ 0 } &
    \cellcolor{blue1}{ 0 } &
    \cellcolor{blue1}{ 0 } &
    \cellcolor{blue1}{0} \\\hline
  \cellcolor{blue1}{ CH Teórica } &
    \cellcolor{blue1}{CH Prática} &
    \cellcolor{blue1}{CH Humanidades} &
    \cellcolor{blue1}{CH Extensão} &
    \cellcolor{blue1}{CH EAD TOTAL} \\\hline
  \cellcolor{blue1}{ 15 } &
    \cellcolor{blue1}{ 15 } &
    \cellcolor{blue1}{ 30 } &
    \cellcolor{blue1}{ 0 } &
    \cellcolor{blue1}{ 0 } \\\hline
  \cellcolor{blue1}{CARGA-HORÁRIA TOTAL} & \multicolumn{4}{p{10cm}|}{ \cellcolor{blue1}{ 30 } } \\\hline
  Área de conhecimento & \multicolumn{4}{p{10cm}|}{ Humanidades } \\\hline
  Ementa & \multicolumn{4}{p{10cm}|}{ Abordagem de gêneros textuais escritos relevantes para a formação acadêmica e profissional: foco em textos aplicados ao cotidiano do Bacharel em Sistemas de Informação. Métodos para produção e interpretação de textos: uso de tecnologias digitais como ferramentas de apoio à comunicação escrita. Características, etapas e mecanismos de construção textual: relação entre a estrutura textual e as demandas da profissão. Avaliação e produção de textos acadêmicos: desenvolvimento de habilidades para avaliar textos acadêmicos e técnicas para construção de textos de forma clara, coerente e coesa. } \\\hline
  Competências & \multicolumn{4}{p{10cm}|}{ C8 } \\\hline


  
	\end{tabular}

\end{quadro}


\clearpage
\newpage

\begin{quadro}[ht!]
  \centering\scriptsize
\caption{Unidade Curricular Introdução à Ciência de Dados e Inteligência Artificial}

\begin{tabular}{|p{2.5cm}|p{2.5cm}|p{2.5cm}|p{2.5cm}|p{2.5cm}|}\hline

  Unidade curricular & \multicolumn{4}{p{10cm}|}{ Introdução à Ciência de Dados e Inteligência Artificial }\\\hline
  Código & 1.4 & Período & \multicolumn{2}{p{5cm}|}{ 1 }\\\hline
  Pré-Requisitos & \multicolumn{4}{p{10cm}|}{  } \\\hline
  \cellcolor{blue1}{ CH PRESENCIAL } &
    \cellcolor{blue1}{CH EAD Síncrona} &
    \cellcolor{blue1}{CH EAD Síncrona Mediada} &
    \cellcolor{blue1}{CH EAD Assíncrona} &
    \cellcolor{blue1}{CH APCC} \\\hline
  \cellcolor{blue1}{ 60 } &
    \cellcolor{blue1}{ 0 } &
    \cellcolor{blue1}{ 0 } &
    \cellcolor{blue1}{ 0 } &
    \cellcolor{blue1}{0} \\\hline
  \cellcolor{blue1}{ CH Teórica } &
    \cellcolor{blue1}{CH Prática} &
    \cellcolor{blue1}{CH Humanidades} &
    \cellcolor{blue1}{CH Extensão} &
    \cellcolor{blue1}{CH EAD TOTAL} \\\hline
  \cellcolor{blue1}{ 30 } &
    \cellcolor{blue1}{ 30 } &
    \cellcolor{blue1}{ 0 } &
    \cellcolor{blue1}{ 0 } &
    \cellcolor{blue1}{ 0 } \\\hline
  \cellcolor{blue1}{CARGA-HORÁRIA TOTAL} & \multicolumn{4}{p{10cm}|}{ \cellcolor{blue1}{ 60 } } \\\hline
  Área de conhecimento & \multicolumn{4}{p{10cm}|}{ Ciência de Dados e Inteligência Artificial } \\\hline
  Ementa & \multicolumn{4}{p{10cm}|}{ Panorama da Ciência de Dados e da Inteligência Artificial. Perfil e áreas de atuação do cientista de dados, competências profissionais e relação com o perfil do egresso da UTFPR. Visão geral do ciclo de projetos em Ciência de Dados e dos ambientes de trabalho. Noções introdutórias de ética, responsabilidade e impactos sociais do uso de dados e IA. Fundamentos básicos de infraestrutura computacional para Ciência de Dados: organização de computadores, sistemas operacionais e uso inicial de ambientes Linux para execução de aplicações e análise de dados. } \\\hline
  Competências & \multicolumn{4}{p{10cm}|}{ C6 } \\\hline


  
	\end{tabular}

\end{quadro}


\clearpage
\newpage

\begin{quadro}[ht!]
  \centering\scriptsize
\caption{Unidade Curricular Matemática Discreta}

\begin{tabular}{|p{2.5cm}|p{2.5cm}|p{2.5cm}|p{2.5cm}|p{2.5cm}|}\hline

  Unidade curricular & \multicolumn{4}{p{10cm}|}{ Matemática Discreta }\\\hline
  Código & 1.5 & Período & \multicolumn{2}{p{5cm}|}{ 1 }\\\hline
  Pré-Requisitos & \multicolumn{4}{p{10cm}|}{  } \\\hline
  \cellcolor{blue1}{ CH PRESENCIAL } &
    \cellcolor{blue1}{CH EAD Síncrona} &
    \cellcolor{blue1}{CH EAD Síncrona Mediada} &
    \cellcolor{blue1}{CH EAD Assíncrona} &
    \cellcolor{blue1}{CH APCC} \\\hline
  \cellcolor{blue1}{ 60 } &
    \cellcolor{blue1}{ 0 } &
    \cellcolor{blue1}{ 0 } &
    \cellcolor{blue1}{ 0 } &
    \cellcolor{blue1}{0} \\\hline
  \cellcolor{blue1}{ CH Teórica } &
    \cellcolor{blue1}{CH Prática} &
    \cellcolor{blue1}{CH Humanidades} &
    \cellcolor{blue1}{CH Extensão} &
    \cellcolor{blue1}{CH EAD TOTAL} \\\hline
  \cellcolor{blue1}{ 30 } &
    \cellcolor{blue1}{ 30 } &
    \cellcolor{blue1}{ 0 } &
    \cellcolor{blue1}{ 0 } &
    \cellcolor{blue1}{ 0 } \\\hline
  \cellcolor{blue1}{CARGA-HORÁRIA TOTAL} & \multicolumn{4}{p{10cm}|}{ \cellcolor{blue1}{ 60 } } \\\hline
  Área de conhecimento & \multicolumn{4}{p{10cm}|}{ Matemática e Estatística } \\\hline
  Ementa & \multicolumn{4}{p{10cm}|}{ Conjuntos Numéricos: tipos e operações com conjuntos. Funções de uma variável: principais funções e propriedades. Lógica Matemática: tabelas verdade, equivalências lógicas, lógica proposicional, indução matemática e raciocínio indutivo. Conceitos básicos de grafos. Representação de grafos. Propriedades e aplicações. Analise Combinatória: princípios fundamentais de contagem. Permutações e combinações. Binômio de Newton. } \\\hline
  Competências & \multicolumn{4}{p{10cm}|}{ C1 } \\\hline


  
	\end{tabular}

\end{quadro}


\clearpage
\newpage

\begin{quadro}[ht!]
  \centering\scriptsize
\caption{Unidade Curricular Projeto Extensionista}

\begin{tabular}{|p{2.5cm}|p{2.5cm}|p{2.5cm}|p{2.5cm}|p{2.5cm}|}\hline

  Unidade curricular & \multicolumn{4}{p{10cm}|}{ Projeto Extensionista }\\\hline
  Código & 1.6 & Período & \multicolumn{2}{p{5cm}|}{ 1 }\\\hline
  Pré-Requisitos & \multicolumn{4}{p{10cm}|}{  } \\\hline
  \cellcolor{blue1}{ CH PRESENCIAL } &
    \cellcolor{blue1}{CH EAD Síncrona} &
    \cellcolor{blue1}{CH EAD Síncrona Mediada} &
    \cellcolor{blue1}{CH EAD Assíncrona} &
    \cellcolor{blue1}{CH APCC} \\\hline
  \cellcolor{blue1}{ 120 } &
    \cellcolor{blue1}{ 0 } &
    \cellcolor{blue1}{ 0 } &
    \cellcolor{blue1}{ 0 } &
    \cellcolor{blue1}{0} \\\hline
  \cellcolor{blue1}{ CH Teórica } &
    \cellcolor{blue1}{CH Prática} &
    \cellcolor{blue1}{CH Humanidades} &
    \cellcolor{blue1}{CH Extensão} &
    \cellcolor{blue1}{CH EAD TOTAL} \\\hline
  \cellcolor{blue1}{ 0 } &
    \cellcolor{blue1}{ 0 } &
    \cellcolor{blue1}{ 0 } &
    \cellcolor{blue1}{ 120 } &
    \cellcolor{blue1}{ 0 } \\\hline
  \cellcolor{blue1}{CARGA-HORÁRIA TOTAL} & \multicolumn{4}{p{10cm}|}{ \cellcolor{blue1}{ 120 } } \\\hline
  Área de conhecimento & \multicolumn{4}{p{10cm}|}{ Extensão } \\\hline
  Ementa & \multicolumn{4}{p{10cm}|}{ Participação em projeto extensionista do campus. } \\\hline
  Competências & \multicolumn{4}{p{10cm}|}{ C6 } \\\hline


  
	\end{tabular}

\end{quadro}


\clearpage
\newpage

\begin{quadro}[ht!]
  \centering\scriptsize
\caption{Unidade Curricular Álgebra Linear}

\begin{tabular}{|p{2.5cm}|p{2.5cm}|p{2.5cm}|p{2.5cm}|p{2.5cm}|}\hline

  Unidade curricular & \multicolumn{4}{p{10cm}|}{ Álgebra Linear }\\\hline
  Código & 2.1 & Período & \multicolumn{2}{p{5cm}|}{ 2 }\\\hline
  Pré-Requisitos & \multicolumn{4}{p{10cm}|}{  } \\\hline
  \cellcolor{blue1}{ CH PRESENCIAL } &
    \cellcolor{blue1}{CH EAD Síncrona} &
    \cellcolor{blue1}{CH EAD Síncrona Mediada} &
    \cellcolor{blue1}{CH EAD Assíncrona} &
    \cellcolor{blue1}{CH APCC} \\\hline
  \cellcolor{blue1}{ 45 } &
    \cellcolor{blue1}{ 0 } &
    \cellcolor{blue1}{ 0 } &
    \cellcolor{blue1}{ 0 } &
    \cellcolor{blue1}{0} \\\hline
  \cellcolor{blue1}{ CH Teórica } &
    \cellcolor{blue1}{CH Prática} &
    \cellcolor{blue1}{CH Humanidades} &
    \cellcolor{blue1}{CH Extensão} &
    \cellcolor{blue1}{CH EAD TOTAL} \\\hline
  \cellcolor{blue1}{ 45 } &
    \cellcolor{blue1}{ 0 } &
    \cellcolor{blue1}{ 0 } &
    \cellcolor{blue1}{ 0 } &
    \cellcolor{blue1}{ 0 } \\\hline
  \cellcolor{blue1}{CARGA-HORÁRIA TOTAL} & \multicolumn{4}{p{10cm}|}{ \cellcolor{blue1}{ 45 } } \\\hline
  Área de conhecimento & \multicolumn{4}{p{10cm}|}{ Matemática e Estatística } \\\hline
  Ementa & \multicolumn{4}{p{10cm}|}{ Conceitos básicos de vetores e matrizes, operações e interpretação geométrica. Sistemas de equações lineares e sua resolução, com aplicações simples. Noções de espaços vetoriais e subespaços, base e dimensão, com ênfase intuitiva. Transformações lineares e representação matricial, com interpretação geométrica. Introdução a autovalores e autovetores, com aplicações em redução de dimensionalidade. Noções de decomposições matriciais (SVD). Normas e distância entre vetores, similaridade e interpretação geométrica.  } \\\hline
  Competências & \multicolumn{4}{p{10cm}|}{ C1 } \\\hline


  
	\end{tabular}

\end{quadro}


\clearpage
\newpage

\begin{quadro}[ht!]
  \centering\scriptsize
\caption{Unidade Curricular Algoritmos 2}

\begin{tabular}{|p{2.5cm}|p{2.5cm}|p{2.5cm}|p{2.5cm}|p{2.5cm}|}\hline

  Unidade curricular & \multicolumn{4}{p{10cm}|}{ Algoritmos 2 }\\\hline
  Código & 2.2 & Período & \multicolumn{2}{p{5cm}|}{ 2 }\\\hline
  Pré-Requisitos & \multicolumn{4}{p{10cm}|}{ Algoritmos 1 } \\\hline
  \cellcolor{blue1}{ CH PRESENCIAL } &
    \cellcolor{blue1}{CH EAD Síncrona} &
    \cellcolor{blue1}{CH EAD Síncrona Mediada} &
    \cellcolor{blue1}{CH EAD Assíncrona} &
    \cellcolor{blue1}{CH APCC} \\\hline
  \cellcolor{blue1}{ 60 } &
    \cellcolor{blue1}{ 0 } &
    \cellcolor{blue1}{ 0 } &
    \cellcolor{blue1}{ 30 } &
    \cellcolor{blue1}{0} \\\hline
  \cellcolor{blue1}{ CH Teórica } &
    \cellcolor{blue1}{CH Prática} &
    \cellcolor{blue1}{CH Humanidades} &
    \cellcolor{blue1}{CH Extensão} &
    \cellcolor{blue1}{CH EAD TOTAL} \\\hline
  \cellcolor{blue1}{ 45 } &
    \cellcolor{blue1}{ 45 } &
    \cellcolor{blue1}{ 0 } &
    \cellcolor{blue1}{ 0 } &
    \cellcolor{blue1}{ 30 } \\\hline
  \cellcolor{blue1}{CARGA-HORÁRIA TOTAL} & \multicolumn{4}{p{10cm}|}{ \cellcolor{blue1}{ 90 } } \\\hline
  Área de conhecimento & \multicolumn{4}{p{10cm}|}{ Computação } \\\hline
  Ementa & \multicolumn{4}{p{10cm}|}{ Técnicas avançadas de programação para resolução de problemas computacionais. Paradigmas de programação, com ênfase em programação funcional e uso de funções de alta ordem para transformação, filtragem e agregação de dados. Abstração e manipulação de dados por referência lógica. Recursividade e iteração, incluindo mapeamento entre abordagens e análise de desempenho. Uso de estruturas de dados eficientes e operações vetorizadas para processamento de coleções de dados. Aplicações dessas técnicas em problemas típicos de processamento, organização e análise de dados em larga escala. } \\\hline
  Competências & \multicolumn{4}{p{10cm}|}{ C1 } \\\hline


  
	\end{tabular}

\end{quadro}


\clearpage
\newpage

\begin{quadro}[ht!]
  \centering\scriptsize
\caption{Unidade Curricular Fundamentos de Inteligência Artificial}

\begin{tabular}{|p{2.5cm}|p{2.5cm}|p{2.5cm}|p{2.5cm}|p{2.5cm}|}\hline

  Unidade curricular & \multicolumn{4}{p{10cm}|}{ Fundamentos de Inteligência Artificial }\\\hline
  Código & 2.3 & Período & \multicolumn{2}{p{5cm}|}{ 2 }\\\hline
  Pré-Requisitos & \multicolumn{4}{p{10cm}|}{  } \\\hline
  \cellcolor{blue1}{ CH PRESENCIAL } &
    \cellcolor{blue1}{CH EAD Síncrona} &
    \cellcolor{blue1}{CH EAD Síncrona Mediada} &
    \cellcolor{blue1}{CH EAD Assíncrona} &
    \cellcolor{blue1}{CH APCC} \\\hline
  \cellcolor{blue1}{ 45 } &
    \cellcolor{blue1}{ 0 } &
    \cellcolor{blue1}{ 0 } &
    \cellcolor{blue1}{ 15 } &
    \cellcolor{blue1}{0} \\\hline
  \cellcolor{blue1}{ CH Teórica } &
    \cellcolor{blue1}{CH Prática} &
    \cellcolor{blue1}{CH Humanidades} &
    \cellcolor{blue1}{CH Extensão} &
    \cellcolor{blue1}{CH EAD TOTAL} \\\hline
  \cellcolor{blue1}{ 30 } &
    \cellcolor{blue1}{ 30 } &
    \cellcolor{blue1}{ 0 } &
    \cellcolor{blue1}{ 0 } &
    \cellcolor{blue1}{ 15 } \\\hline
  \cellcolor{blue1}{CARGA-HORÁRIA TOTAL} & \multicolumn{4}{p{10cm}|}{ \cellcolor{blue1}{ 60 } } \\\hline
  Área de conhecimento & \multicolumn{4}{p{10cm}|}{ Ciência de Dados e Inteligência Artificial } \\\hline
  Ementa & \multicolumn{4}{p{10cm}|}{ Introdução à Inteligência Artificial: conceitos básicos, histórico e evolução da área. Principais abordagens e taxonomias de IA. Noções de agentes inteligentes, percepção, decisão e raciocínio. Visão geral dos principais paradigmas e aplicações da IA. Tendências atuais, IA generativa e modelos de grande escala. Aspectos éticos, sociais e profissionais relacionados ao uso da Inteligência Artificial. } \\\hline
  Competências & \multicolumn{4}{p{10cm}|}{ C4 } \\\hline


  
	\end{tabular}

\end{quadro}


\clearpage
\newpage

\begin{quadro}[ht!]
  \centering\scriptsize
\caption{Unidade Curricular Gestão Empresarial}

\begin{tabular}{|p{2.5cm}|p{2.5cm}|p{2.5cm}|p{2.5cm}|p{2.5cm}|}\hline

  Unidade curricular & \multicolumn{4}{p{10cm}|}{ Gestão Empresarial }\\\hline
  Código & 2.4 & Período & \multicolumn{2}{p{5cm}|}{ 2 }\\\hline
  Pré-Requisitos & \multicolumn{4}{p{10cm}|}{  } \\\hline
  \cellcolor{blue1}{ CH PRESENCIAL } &
    \cellcolor{blue1}{CH EAD Síncrona} &
    \cellcolor{blue1}{CH EAD Síncrona Mediada} &
    \cellcolor{blue1}{CH EAD Assíncrona} &
    \cellcolor{blue1}{CH APCC} \\\hline
  \cellcolor{blue1}{ 0 } &
    \cellcolor{blue1}{ 0 } &
    \cellcolor{blue1}{ 30 } &
    \cellcolor{blue1}{ 0 } &
    \cellcolor{blue1}{0} \\\hline
  \cellcolor{blue1}{ CH Teórica } &
    \cellcolor{blue1}{CH Prática} &
    \cellcolor{blue1}{CH Humanidades} &
    \cellcolor{blue1}{CH Extensão} &
    \cellcolor{blue1}{CH EAD TOTAL} \\\hline
  \cellcolor{blue1}{ 15 } &
    \cellcolor{blue1}{ 15 } &
    \cellcolor{blue1}{ 30 } &
    \cellcolor{blue1}{ 0 } &
    \cellcolor{blue1}{ 30 } \\\hline
  \cellcolor{blue1}{CARGA-HORÁRIA TOTAL} & \multicolumn{4}{p{10cm}|}{ \cellcolor{blue1}{ 30 } } \\\hline
  Área de conhecimento & \multicolumn{4}{p{10cm}|}{ Humanidades } \\\hline
  Ementa & \multicolumn{4}{p{10cm}|}{ Aborda a administração de empresas no contexto de origens, funções e a atuação do administrador. Nela, os estudantes aprendem os conceitos básicos de administração e toma contato com ferramentas elementares de planejamento e gerenciamento de empresas. Ao final da unidade curricular os alunos são capazes de utilizar as ferramentas básicas de gestão para tomada de decisão em empresas. } \\\hline
  Competências & \multicolumn{4}{p{10cm}|}{ C7 } \\\hline


  
	\end{tabular}

\end{quadro}


\clearpage
\newpage

\begin{quadro}[ht!]
  \centering\scriptsize
\caption{Unidade Curricular Probabilidade e Estatística}

\begin{tabular}{|p{2.5cm}|p{2.5cm}|p{2.5cm}|p{2.5cm}|p{2.5cm}|}\hline

  Unidade curricular & \multicolumn{4}{p{10cm}|}{ Probabilidade e Estatística }\\\hline
  Código & 2.5 & Período & \multicolumn{2}{p{5cm}|}{ 2 }\\\hline
  Pré-Requisitos & \multicolumn{4}{p{10cm}|}{  } \\\hline
  \cellcolor{blue1}{ CH PRESENCIAL } &
    \cellcolor{blue1}{CH EAD Síncrona} &
    \cellcolor{blue1}{CH EAD Síncrona Mediada} &
    \cellcolor{blue1}{CH EAD Assíncrona} &
    \cellcolor{blue1}{CH APCC} \\\hline
  \cellcolor{blue1}{ 60 } &
    \cellcolor{blue1}{ 0 } &
    \cellcolor{blue1}{ 0 } &
    \cellcolor{blue1}{ 0 } &
    \cellcolor{blue1}{0} \\\hline
  \cellcolor{blue1}{ CH Teórica } &
    \cellcolor{blue1}{CH Prática} &
    \cellcolor{blue1}{CH Humanidades} &
    \cellcolor{blue1}{CH Extensão} &
    \cellcolor{blue1}{CH EAD TOTAL} \\\hline
  \cellcolor{blue1}{ 30 } &
    \cellcolor{blue1}{ 30 } &
    \cellcolor{blue1}{ 0 } &
    \cellcolor{blue1}{ 0 } &
    \cellcolor{blue1}{ 0 } \\\hline
  \cellcolor{blue1}{CARGA-HORÁRIA TOTAL} & \multicolumn{4}{p{10cm}|}{ \cellcolor{blue1}{ 60 } } \\\hline
  Área de conhecimento & \multicolumn{4}{p{10cm}|}{ Matemática e Estatística } \\\hline
  Ementa & \multicolumn{4}{p{10cm}|}{ Estatística descritiva: medidas de tendência central, medidas de dispersão e representações gráficas. Noções de amostragem. Introdução a probabilidade: probabilidade clássica e condicional. Distribuição discreta de probabilidades: binomial. Distribuição continua de probabilidade: normal. Estimação de parâmetros: pontual e intervalar. Teste de hipótese para média: uma amostra, duas amostras independentes e duas amostras dependentes. Teste de normalidade. } \\\hline
  Competências & \multicolumn{4}{p{10cm}|}{ C1 } \\\hline


  
	\end{tabular}

\end{quadro}


\clearpage
\newpage

\begin{quadro}[ht!]
  \centering\scriptsize
\caption{Unidade Curricular Programação Orientada a Objetos}

\begin{tabular}{|p{2.5cm}|p{2.5cm}|p{2.5cm}|p{2.5cm}|p{2.5cm}|}\hline

  Unidade curricular & \multicolumn{4}{p{10cm}|}{ Programação Orientada a Objetos }\\\hline
  Código & 2.6 & Período & \multicolumn{2}{p{5cm}|}{ 2 }\\\hline
  Pré-Requisitos & \multicolumn{4}{p{10cm}|}{ Algoritmos 1 } \\\hline
  \cellcolor{blue1}{ CH PRESENCIAL } &
    \cellcolor{blue1}{CH EAD Síncrona} &
    \cellcolor{blue1}{CH EAD Síncrona Mediada} &
    \cellcolor{blue1}{CH EAD Assíncrona} &
    \cellcolor{blue1}{CH APCC} \\\hline
  \cellcolor{blue1}{ 30 } &
    \cellcolor{blue1}{ 0 } &
    \cellcolor{blue1}{ 0 } &
    \cellcolor{blue1}{ 30 } &
    \cellcolor{blue1}{0} \\\hline
  \cellcolor{blue1}{ CH Teórica } &
    \cellcolor{blue1}{CH Prática} &
    \cellcolor{blue1}{CH Humanidades} &
    \cellcolor{blue1}{CH Extensão} &
    \cellcolor{blue1}{CH EAD TOTAL} \\\hline
  \cellcolor{blue1}{ 30 } &
    \cellcolor{blue1}{ 30 } &
    \cellcolor{blue1}{ 0 } &
    \cellcolor{blue1}{ 0 } &
    \cellcolor{blue1}{ 30 } \\\hline
  \cellcolor{blue1}{CARGA-HORÁRIA TOTAL} & \multicolumn{4}{p{10cm}|}{ \cellcolor{blue1}{ 60 } } \\\hline
  Área de conhecimento & \multicolumn{4}{p{10cm}|}{ Computação } \\\hline
  Ementa & \multicolumn{4}{p{10cm}|}{ Fundamentos de orientação a objetos: definições, artefatos, atributos, métodos, objetos, tipos de classes, operador de casting e o paralelo entre orientação à objetos e o paradigma estruturado. Relacionamento entre classes e objetos: encapsulamento, herança, tipos de polimorfismo e interfaces. Tratamento de exceções: definição, mecanismos de disparo e tratamento e classes de exceção. Desenvolvimento de aplicações: utilização de IDEs para implementação de aplicações utilizando paradigma de orientação à objetos. } \\\hline
  Competências & \multicolumn{4}{p{10cm}|}{ C1 } \\\hline


  
	\end{tabular}

\end{quadro}


\clearpage
\newpage

\begin{quadro}[ht!]
  \centering\scriptsize
\caption{Unidade Curricular Análise Estatística}

\begin{tabular}{|p{2.5cm}|p{2.5cm}|p{2.5cm}|p{2.5cm}|p{2.5cm}|}\hline

  Unidade curricular & \multicolumn{4}{p{10cm}|}{ Análise Estatística }\\\hline
  Código & 3.1 & Período & \multicolumn{2}{p{5cm}|}{ 3 }\\\hline
  Pré-Requisitos & \multicolumn{4}{p{10cm}|}{ Probabilidade e Estatística } \\\hline
  \cellcolor{blue1}{ CH PRESENCIAL } &
    \cellcolor{blue1}{CH EAD Síncrona} &
    \cellcolor{blue1}{CH EAD Síncrona Mediada} &
    \cellcolor{blue1}{CH EAD Assíncrona} &
    \cellcolor{blue1}{CH APCC} \\\hline
  \cellcolor{blue1}{ 60 } &
    \cellcolor{blue1}{ 0 } &
    \cellcolor{blue1}{ 0 } &
    \cellcolor{blue1}{ 0 } &
    \cellcolor{blue1}{0} \\\hline
  \cellcolor{blue1}{ CH Teórica } &
    \cellcolor{blue1}{CH Prática} &
    \cellcolor{blue1}{CH Humanidades} &
    \cellcolor{blue1}{CH Extensão} &
    \cellcolor{blue1}{CH EAD TOTAL} \\\hline
  \cellcolor{blue1}{ 30 } &
    \cellcolor{blue1}{ 30 } &
    \cellcolor{blue1}{ 0 } &
    \cellcolor{blue1}{ 0 } &
    \cellcolor{blue1}{ 0 } \\\hline
  \cellcolor{blue1}{CARGA-HORÁRIA TOTAL} & \multicolumn{4}{p{10cm}|}{ \cellcolor{blue1}{ 60 } } \\\hline
  Área de conhecimento & \multicolumn{4}{p{10cm}|}{ Matemática e Estatística } \\\hline
  Ementa & \multicolumn{4}{p{10cm}|}{ Testes não paramétricas: uma amostra (sign test), duas amostras independente, duas amostras pareadas (Wilcoxon signed-rank). Teste de aderência de distribuição. Correlação entre variáveis numéricas. Introdução a regressão e modelos lineares/generalizados. Análise de resíduos. ANOVA: um fator e dois ou mais fatores. Pressupostos da ANOVA (teste de normalidade e homogeneidade de variâncias). Tamanhos de efeitos. Teste de comparação múltipla. ANOVA com medidas repetidas (intra-sujeitos).  Múltiplos grupos independentes (kruskal Wallis). } \\\hline
  Competências & \multicolumn{4}{p{10cm}|}{ C1 } \\\hline


  
	\end{tabular}

\end{quadro}


\clearpage
\newpage

\begin{quadro}[ht!]
  \centering\scriptsize
\caption{Unidade Curricular Bancos de Dados Relacionais}

\begin{tabular}{|p{2.5cm}|p{2.5cm}|p{2.5cm}|p{2.5cm}|p{2.5cm}|}\hline

  Unidade curricular & \multicolumn{4}{p{10cm}|}{ Bancos de Dados Relacionais }\\\hline
  Código & 3.2 & Período & \multicolumn{2}{p{5cm}|}{ 3 }\\\hline
  Pré-Requisitos & \multicolumn{4}{p{10cm}|}{  } \\\hline
  \cellcolor{blue1}{ CH PRESENCIAL } &
    \cellcolor{blue1}{CH EAD Síncrona} &
    \cellcolor{blue1}{CH EAD Síncrona Mediada} &
    \cellcolor{blue1}{CH EAD Assíncrona} &
    \cellcolor{blue1}{CH APCC} \\\hline
  \cellcolor{blue1}{ 60 } &
    \cellcolor{blue1}{ 0 } &
    \cellcolor{blue1}{ 0 } &
    \cellcolor{blue1}{ 0 } &
    \cellcolor{blue1}{0} \\\hline
  \cellcolor{blue1}{ CH Teórica } &
    \cellcolor{blue1}{CH Prática} &
    \cellcolor{blue1}{CH Humanidades} &
    \cellcolor{blue1}{CH Extensão} &
    \cellcolor{blue1}{CH EAD TOTAL} \\\hline
  \cellcolor{blue1}{ 30 } &
    \cellcolor{blue1}{ 30 } &
    \cellcolor{blue1}{ 0 } &
    \cellcolor{blue1}{ 0 } &
    \cellcolor{blue1}{ 0 } \\\hline
  \cellcolor{blue1}{CARGA-HORÁRIA TOTAL} & \multicolumn{4}{p{10cm}|}{ \cellcolor{blue1}{ 60 } } \\\hline
  Área de conhecimento & \multicolumn{4}{p{10cm}|}{ Computação } \\\hline
  Ementa & \multicolumn{4}{p{10cm}|}{ Princípios de banco de dados: conceitos, aplicações e modelo relacional (relação, tupla, atributo, restrições de integridade e normalização) e construção de representações. Linguagem de Definição de Dados (DDL): criação, atualização e regras de integridade. Linguagem de Manipulação de Dados (DML): consultas simples, consultas aninhadas, operações de conjuntos, funções de agregação, inserção, remoção e atualização. Projeto de Banco de Dados: projeto conceitual, projeto lógico e projeto físico. } \\\hline
  Competências & \multicolumn{4}{p{10cm}|}{ C3 } \\\hline


  
	\end{tabular}

\end{quadro}


\clearpage
\newpage

\begin{quadro}[ht!]
  \centering\scriptsize
\caption{Unidade Curricular Engenharia de Dados}

\begin{tabular}{|p{2.5cm}|p{2.5cm}|p{2.5cm}|p{2.5cm}|p{2.5cm}|}\hline

  Unidade curricular & \multicolumn{4}{p{10cm}|}{ Engenharia de Dados }\\\hline
  Código & 3.3 & Período & \multicolumn{2}{p{5cm}|}{ 3 }\\\hline
  Pré-Requisitos & \multicolumn{4}{p{10cm}|}{ Algoritmos 2 } \\\hline
  \cellcolor{blue1}{ CH PRESENCIAL } &
    \cellcolor{blue1}{CH EAD Síncrona} &
    \cellcolor{blue1}{CH EAD Síncrona Mediada} &
    \cellcolor{blue1}{CH EAD Assíncrona} &
    \cellcolor{blue1}{CH APCC} \\\hline
  \cellcolor{blue1}{ 60 } &
    \cellcolor{blue1}{ 0 } &
    \cellcolor{blue1}{ 0 } &
    \cellcolor{blue1}{ 0 } &
    \cellcolor{blue1}{0} \\\hline
  \cellcolor{blue1}{ CH Teórica } &
    \cellcolor{blue1}{CH Prática} &
    \cellcolor{blue1}{CH Humanidades} &
    \cellcolor{blue1}{CH Extensão} &
    \cellcolor{blue1}{CH EAD TOTAL} \\\hline
  \cellcolor{blue1}{ 30 } &
    \cellcolor{blue1}{ 30 } &
    \cellcolor{blue1}{ 0 } &
    \cellcolor{blue1}{ 0 } &
    \cellcolor{blue1}{ 0 } \\\hline
  \cellcolor{blue1}{CARGA-HORÁRIA TOTAL} & \multicolumn{4}{p{10cm}|}{ \cellcolor{blue1}{ 60 } } \\\hline
  Área de conhecimento & \multicolumn{4}{p{10cm}|}{ Ciência de Dados e Inteligência Artificial } \\\hline
  Ementa & \multicolumn{4}{p{10cm}|}{ Fundamentos de Engenharia de Dados e ciclo de vida dos dados em sistemas analíticos. Ingestão, limpeza, armazenamento, padronização, qualidade (curadoria) e transformação de dados. Seleção de características e balanceamento de dados no preparo para análise. Construção de pipelines de dados reprodutíveis, automatizados e monitoráveis, com versionamento e rastreabilidade. Visão de indústria: ETL e ELT, modelagem analítica, data lakes e data warehouses. } \\\hline
  Competências & \multicolumn{4}{p{10cm}|}{ C3 } \\\hline


  
	\end{tabular}

\end{quadro}


\clearpage
\newpage

\begin{quadro}[ht!]
  \centering\scriptsize
\caption{Unidade Curricular Estrutura de Dados 1}

\begin{tabular}{|p{2.5cm}|p{2.5cm}|p{2.5cm}|p{2.5cm}|p{2.5cm}|}\hline

  Unidade curricular & \multicolumn{4}{p{10cm}|}{ Estrutura de Dados 1 }\\\hline
  Código & 3.4 & Período & \multicolumn{2}{p{5cm}|}{ 3 }\\\hline
  Pré-Requisitos & \multicolumn{4}{p{10cm}|}{ Algoritmos 2 } \\\hline
  \cellcolor{blue1}{ CH PRESENCIAL } &
    \cellcolor{blue1}{CH EAD Síncrona} &
    \cellcolor{blue1}{CH EAD Síncrona Mediada} &
    \cellcolor{blue1}{CH EAD Assíncrona} &
    \cellcolor{blue1}{CH APCC} \\\hline
  \cellcolor{blue1}{ 60 } &
    \cellcolor{blue1}{ 0 } &
    \cellcolor{blue1}{ 0 } &
    \cellcolor{blue1}{ 0 } &
    \cellcolor{blue1}{0} \\\hline
  \cellcolor{blue1}{ CH Teórica } &
    \cellcolor{blue1}{CH Prática} &
    \cellcolor{blue1}{CH Humanidades} &
    \cellcolor{blue1}{CH Extensão} &
    \cellcolor{blue1}{CH EAD TOTAL} \\\hline
  \cellcolor{blue1}{ 30 } &
    \cellcolor{blue1}{ 30 } &
    \cellcolor{blue1}{ 0 } &
    \cellcolor{blue1}{ 0 } &
    \cellcolor{blue1}{ 0 } \\\hline
  \cellcolor{blue1}{CARGA-HORÁRIA TOTAL} & \multicolumn{4}{p{10cm}|}{ \cellcolor{blue1}{ 60 } } \\\hline
  Área de conhecimento & \multicolumn{4}{p{10cm}|}{ Computação } \\\hline
  Ementa & \multicolumn{4}{p{10cm}|}{ Tipos Abstratos de Dados (TAD): conceitos, definições e estratégias de implementação. Estruturas lineares: pilhas e filas — operações fundamentais, implementações e aplicações. Listas lineares: definição, tipos, operações, implementação e aplicações. Estruturas hierárquicas: árvores binárias — conceitos, operações, busca e aplicações. Tabelas Hash: princípios, funções de espalhamento, tratamento de colisões e aplicações. } \\\hline
  Competências & \multicolumn{4}{p{10cm}|}{ C1 } \\\hline


  
	\end{tabular}

\end{quadro}


\clearpage
\newpage

\begin{quadro}[ht!]
  \centering\scriptsize
\caption{Unidade Curricular Gerenciamento de Projetos e Desenvolvimento Ágil}

\begin{tabular}{|p{2.5cm}|p{2.5cm}|p{2.5cm}|p{2.5cm}|p{2.5cm}|}\hline

  Unidade curricular & \multicolumn{4}{p{10cm}|}{ Gerenciamento de Projetos e Desenvolvimento Ágil }\\\hline
  Código & 3.5 & Período & \multicolumn{2}{p{5cm}|}{ 3 }\\\hline
  Pré-Requisitos & \multicolumn{4}{p{10cm}|}{  } \\\hline
  \cellcolor{blue1}{ CH PRESENCIAL } &
    \cellcolor{blue1}{CH EAD Síncrona} &
    \cellcolor{blue1}{CH EAD Síncrona Mediada} &
    \cellcolor{blue1}{CH EAD Assíncrona} &
    \cellcolor{blue1}{CH APCC} \\\hline
  \cellcolor{blue1}{ 0 } &
    \cellcolor{blue1}{ 0 } &
    \cellcolor{blue1}{ 30 } &
    \cellcolor{blue1}{ 30 } &
    \cellcolor{blue1}{0} \\\hline
  \cellcolor{blue1}{ CH Teórica } &
    \cellcolor{blue1}{CH Prática} &
    \cellcolor{blue1}{CH Humanidades} &
    \cellcolor{blue1}{CH Extensão} &
    \cellcolor{blue1}{CH EAD TOTAL} \\\hline
  \cellcolor{blue1}{ 30 } &
    \cellcolor{blue1}{ 30 } &
    \cellcolor{blue1}{ 60 } &
    \cellcolor{blue1}{ 0 } &
    \cellcolor{blue1}{ 60 } \\\hline
  \cellcolor{blue1}{CARGA-HORÁRIA TOTAL} & \multicolumn{4}{p{10cm}|}{ \cellcolor{blue1}{ 60 } } \\\hline
  Área de conhecimento & \multicolumn{4}{p{10cm}|}{ Humanidades } \\\hline
  Ementa & \multicolumn{4}{p{10cm}|}{ Conceito e os objetivos da gerência de projetos. Perfil do Gerente de Projetos. Metodologias, técnicas e ferramentas da gerência de projetos. Abertura e definição do escopo de um projeto. Planejamento de um projeto. Execução, acompanhamento e controle de um projeto. Revisão e avaliação de um projeto. Fechamento de um projeto. Metodologias e práticas de gerência de projetos com foco em gestão de projetos de Tecnologia da Informação e Comunicação. Gerência Ágil de Projeto. Práticas de Desenvolvimento Ágil. Metodologias ágeis (Scrum e XP), Engenharia de Requisitos, Modelagem, Qualidade de Software e Gerência de Configuração. } \\\hline
  Competências & \multicolumn{4}{p{10cm}|}{ C7 } \\\hline


  
	\end{tabular}

\end{quadro}


\clearpage
\newpage

\begin{quadro}[ht!]
  \centering\scriptsize
\caption{Unidade Curricular Aprendizado de Máquina 1}

\begin{tabular}{|p{2.5cm}|p{2.5cm}|p{2.5cm}|p{2.5cm}|p{2.5cm}|}\hline

  Unidade curricular & \multicolumn{4}{p{10cm}|}{ Aprendizado de Máquina 1 }\\\hline
  Código & 4.1 & Período & \multicolumn{2}{p{5cm}|}{ 4 }\\\hline
  Pré-Requisitos & \multicolumn{4}{p{10cm}|}{ Algoritmos 2 } \\\hline
  \cellcolor{blue1}{ CH PRESENCIAL } &
    \cellcolor{blue1}{CH EAD Síncrona} &
    \cellcolor{blue1}{CH EAD Síncrona Mediada} &
    \cellcolor{blue1}{CH EAD Assíncrona} &
    \cellcolor{blue1}{CH APCC} \\\hline
  \cellcolor{blue1}{ 60 } &
    \cellcolor{blue1}{ 0 } &
    \cellcolor{blue1}{ 0 } &
    \cellcolor{blue1}{ 30 } &
    \cellcolor{blue1}{0} \\\hline
  \cellcolor{blue1}{ CH Teórica } &
    \cellcolor{blue1}{CH Prática} &
    \cellcolor{blue1}{CH Humanidades} &
    \cellcolor{blue1}{CH Extensão} &
    \cellcolor{blue1}{CH EAD TOTAL} \\\hline
  \cellcolor{blue1}{ 30 } &
    \cellcolor{blue1}{ 30 } &
    \cellcolor{blue1}{ 0 } &
    \cellcolor{blue1}{ 0 } &
    \cellcolor{blue1}{ 30 } \\\hline
  \cellcolor{blue1}{CARGA-HORÁRIA TOTAL} & \multicolumn{4}{p{10cm}|}{ \cellcolor{blue1}{ 90 } } \\\hline
  Área de conhecimento & \multicolumn{4}{p{10cm}|}{ Ciência de Dados e Inteligência Artificial } \\\hline
  Ementa & \multicolumn{4}{p{10cm}|}{ Fundamentos de aprendizado de máquina: tipos de aprendizado, tarefas e etapas, probabilidade condicional, espaço amostral, teoria de decisão Bayesiana, vetores e espaços de características. Aprendizado não supervisionado: métricas de distância e similaridade, algoritmos de clustering baseada em distancia, agrupamento hierarquico. Aprendizado supervisionado: modelos preditivos (Árvore de Decisão, Redes Bayesianas Naive, Redes de Crença, Regressão Logística, SVM), agrupamento de dados, metodologia de validação e avaliação, aplicações e tendências. Metodologia de experimentação: métricas de avaliação (Acurácia, Precisão, Revocação, F-Score, ROC, ROC-AUC), métodos de validação (k-fold, holdout, bootstrap, leave-one-out), Overfitting e Underrfitting, análise de resultados. } \\\hline
  Competências & \multicolumn{4}{p{10cm}|}{ C4 } \\\hline


  
	\end{tabular}

\end{quadro}


\clearpage
\newpage

\begin{quadro}[ht!]
  \centering\scriptsize
\caption{Unidade Curricular Bancos de Dados Não Relacionais}

\begin{tabular}{|p{2.5cm}|p{2.5cm}|p{2.5cm}|p{2.5cm}|p{2.5cm}|}\hline

  Unidade curricular & \multicolumn{4}{p{10cm}|}{ Bancos de Dados Não Relacionais }\\\hline
  Código & 4.2 & Período & \multicolumn{2}{p{5cm}|}{ 4 }\\\hline
  Pré-Requisitos & \multicolumn{4}{p{10cm}|}{  } \\\hline
  \cellcolor{blue1}{ CH PRESENCIAL } &
    \cellcolor{blue1}{CH EAD Síncrona} &
    \cellcolor{blue1}{CH EAD Síncrona Mediada} &
    \cellcolor{blue1}{CH EAD Assíncrona} &
    \cellcolor{blue1}{CH APCC} \\\hline
  \cellcolor{blue1}{ 30 } &
    \cellcolor{blue1}{ 0 } &
    \cellcolor{blue1}{ 0 } &
    \cellcolor{blue1}{ 30 } &
    \cellcolor{blue1}{0} \\\hline
  \cellcolor{blue1}{ CH Teórica } &
    \cellcolor{blue1}{CH Prática} &
    \cellcolor{blue1}{CH Humanidades} &
    \cellcolor{blue1}{CH Extensão} &
    \cellcolor{blue1}{CH EAD TOTAL} \\\hline
  \cellcolor{blue1}{ 30 } &
    \cellcolor{blue1}{ 30 } &
    \cellcolor{blue1}{ 0 } &
    \cellcolor{blue1}{ 0 } &
    \cellcolor{blue1}{ 30 } \\\hline
  \cellcolor{blue1}{CARGA-HORÁRIA TOTAL} & \multicolumn{4}{p{10cm}|}{ \cellcolor{blue1}{ 60 } } \\\hline
  Área de conhecimento & \multicolumn{4}{p{10cm}|}{ Ciência de Dados e Inteligência Artificial } \\\hline
  Ementa & \multicolumn{4}{p{10cm}|}{ Fundamentos de Banco de Dados NoSQL. Modelos de dados. Gerenciamento de dados. Otimização específica para aplicativos que manipulam grande volume de dados. Entendimento e aplicação do conceito de baixa latência, Flexibilidade, Escalabilidade, Alta performance, Funcionalidade. Tipos de registro de tabelas no banco de dados NoSQL. } \\\hline
  Competências & \multicolumn{4}{p{10cm}|}{ C3 } \\\hline


  
	\end{tabular}

\end{quadro}


\clearpage
\newpage

\begin{quadro}[ht!]
  \centering\scriptsize
\caption{Unidade Curricular Cálculo Numérico}

\begin{tabular}{|p{2.5cm}|p{2.5cm}|p{2.5cm}|p{2.5cm}|p{2.5cm}|}\hline

  Unidade curricular & \multicolumn{4}{p{10cm}|}{ Cálculo Numérico }\\\hline
  Código & 4.3 & Período & \multicolumn{2}{p{5cm}|}{ 4 }\\\hline
  Pré-Requisitos & \multicolumn{4}{p{10cm}|}{ Algoritmos 2 } \\\hline
  \cellcolor{blue1}{ CH PRESENCIAL } &
    \cellcolor{blue1}{CH EAD Síncrona} &
    \cellcolor{blue1}{CH EAD Síncrona Mediada} &
    \cellcolor{blue1}{CH EAD Assíncrona} &
    \cellcolor{blue1}{CH APCC} \\\hline
  \cellcolor{blue1}{ 60 } &
    \cellcolor{blue1}{ 0 } &
    \cellcolor{blue1}{ 0 } &
    \cellcolor{blue1}{ 0 } &
    \cellcolor{blue1}{0} \\\hline
  \cellcolor{blue1}{ CH Teórica } &
    \cellcolor{blue1}{CH Prática} &
    \cellcolor{blue1}{CH Humanidades} &
    \cellcolor{blue1}{CH Extensão} &
    \cellcolor{blue1}{CH EAD TOTAL} \\\hline
  \cellcolor{blue1}{ 60 } &
    \cellcolor{blue1}{ 0 } &
    \cellcolor{blue1}{ 0 } &
    \cellcolor{blue1}{ 0 } &
    \cellcolor{blue1}{ 0 } \\\hline
  \cellcolor{blue1}{CARGA-HORÁRIA TOTAL} & \multicolumn{4}{p{10cm}|}{ \cellcolor{blue1}{ 60 } } \\\hline
  Área de conhecimento & \multicolumn{4}{p{10cm}|}{ Matemática e Estatística } \\\hline
  Ementa & \multicolumn{4}{p{10cm}|}{ Funções reais de variáveis, limites, continuidade e derivadas. Otimização de funções de gradiente, máximos e mínimos, convexidade e problemas sem e com restrições. Integrais definidas e noções de integrais múltiplas e aplicações. Introdução à otimização numérica considerando o método do gradiente descendente e sua relação com o treinamento de modelos de aprendizado de máquina. } \\\hline
  Competências & \multicolumn{4}{p{10cm}|}{ C1 } \\\hline


  
	\end{tabular}

\end{quadro}


\clearpage
\newpage

\begin{quadro}[ht!]
  \centering\scriptsize
\caption{Unidade Curricular Estrutura de Dados 2}

\begin{tabular}{|p{2.5cm}|p{2.5cm}|p{2.5cm}|p{2.5cm}|p{2.5cm}|}\hline

  Unidade curricular & \multicolumn{4}{p{10cm}|}{ Estrutura de Dados 2 }\\\hline
  Código & 4.4 & Período & \multicolumn{2}{p{5cm}|}{ 4 }\\\hline
  Pré-Requisitos & \multicolumn{4}{p{10cm}|}{ Estrutura de Dados 1 } \\\hline
  \cellcolor{blue1}{ CH PRESENCIAL } &
    \cellcolor{blue1}{CH EAD Síncrona} &
    \cellcolor{blue1}{CH EAD Síncrona Mediada} &
    \cellcolor{blue1}{CH EAD Assíncrona} &
    \cellcolor{blue1}{CH APCC} \\\hline
  \cellcolor{blue1}{ 60 } &
    \cellcolor{blue1}{ 0 } &
    \cellcolor{blue1}{ 0 } &
    \cellcolor{blue1}{ 0 } &
    \cellcolor{blue1}{0} \\\hline
  \cellcolor{blue1}{ CH Teórica } &
    \cellcolor{blue1}{CH Prática} &
    \cellcolor{blue1}{CH Humanidades} &
    \cellcolor{blue1}{CH Extensão} &
    \cellcolor{blue1}{CH EAD TOTAL} \\\hline
  \cellcolor{blue1}{ 30 } &
    \cellcolor{blue1}{ 30 } &
    \cellcolor{blue1}{ 0 } &
    \cellcolor{blue1}{ 0 } &
    \cellcolor{blue1}{ 0 } \\\hline
  \cellcolor{blue1}{CARGA-HORÁRIA TOTAL} & \multicolumn{4}{p{10cm}|}{ \cellcolor{blue1}{ 60 } } \\\hline
  Área de conhecimento & \multicolumn{4}{p{10cm}|}{ Computação } \\\hline
  Ementa & \multicolumn{4}{p{10cm}|}{ Algoritmos de ordenação: princípios, classificação, funcionamento e aplicações. Análise introdutória de complexidade de algoritmos, com foco em tempo e espaço. Árvores balanceadas: conceitos, operações fundamentais, busca e aplicações, incluindo árvores AVL, rubro-negras e noções de árvores B. Introdução a grafos: conceitos básicos, formas de representação, operações e algoritmos de busca, com aplicações em problemas computacionais. } \\\hline
  Competências & \multicolumn{4}{p{10cm}|}{ C1 } \\\hline


  
	\end{tabular}

\end{quadro}


\clearpage
\newpage

\begin{quadro}[ht!]
  \centering\scriptsize
\caption{Unidade Curricular Sistemas de Computação e Comunicação de Dados}

\begin{tabular}{|p{2.5cm}|p{2.5cm}|p{2.5cm}|p{2.5cm}|p{2.5cm}|}\hline

  Unidade curricular & \multicolumn{4}{p{10cm}|}{ Sistemas de Computação e Comunicação de Dados }\\\hline
  Código & 4.5 & Período & \multicolumn{2}{p{5cm}|}{ 4 }\\\hline
  Pré-Requisitos & \multicolumn{4}{p{10cm}|}{  } \\\hline
  \cellcolor{blue1}{ CH PRESENCIAL } &
    \cellcolor{blue1}{CH EAD Síncrona} &
    \cellcolor{blue1}{CH EAD Síncrona Mediada} &
    \cellcolor{blue1}{CH EAD Assíncrona} &
    \cellcolor{blue1}{CH APCC} \\\hline
  \cellcolor{blue1}{ 30 } &
    \cellcolor{blue1}{ 0 } &
    \cellcolor{blue1}{ 0 } &
    \cellcolor{blue1}{ 30 } &
    \cellcolor{blue1}{0} \\\hline
  \cellcolor{blue1}{ CH Teórica } &
    \cellcolor{blue1}{CH Prática} &
    \cellcolor{blue1}{CH Humanidades} &
    \cellcolor{blue1}{CH Extensão} &
    \cellcolor{blue1}{CH EAD TOTAL} \\\hline
  \cellcolor{blue1}{ 30 } &
    \cellcolor{blue1}{ 30 } &
    \cellcolor{blue1}{ 0 } &
    \cellcolor{blue1}{ 0 } &
    \cellcolor{blue1}{ 30 } \\\hline
  \cellcolor{blue1}{CARGA-HORÁRIA TOTAL} & \multicolumn{4}{p{10cm}|}{ \cellcolor{blue1}{ 60 } } \\\hline
  Área de conhecimento & \multicolumn{4}{p{10cm}|}{ Ciência de Dados e Inteligência Artificial } \\\hline
  Ementa & \multicolumn{4}{p{10cm}|}{ Sistemas operacionais: classificação, estrutura, componentes, concorrência, configuração e utilização de sistemas Linux. Processos: definição e exemplos dos uso de threads, escalonamento, sincronização e deadlocks. Memória: formas de endereçamento, gerenciamento de memória virtual e real, alocação, swapping, segmentação e paginação de memória. Gerenciamento de dispositivos de entrada e saída: discos, interfaces com o usuário e o conceito de spooling. Princípios de Redes de Computadores e a Internet: conceitos, definições e aplicações, arquitetura de redes, topologia, arquitetura em camadas da Internet. Camada de Aplicação: conceitos e utilização dos protocolos HTTP, SSH, FTP, SMTP, IMAP, POP3, DNS, P2P, VPN. } \\\hline
  Competências & \multicolumn{4}{p{10cm}|}{ C5 } \\\hline


  
	\end{tabular}

\end{quadro}


\clearpage
\newpage

\begin{quadro}[ht!]
  \centering\scriptsize
\caption{Unidade Curricular Aprendizado de Máquina 2}

\begin{tabular}{|p{2.5cm}|p{2.5cm}|p{2.5cm}|p{2.5cm}|p{2.5cm}|}\hline

  Unidade curricular & \multicolumn{4}{p{10cm}|}{ Aprendizado de Máquina 2 }\\\hline
  Código & 5.1 & Período & \multicolumn{2}{p{5cm}|}{ 5 }\\\hline
  Pré-Requisitos & \multicolumn{4}{p{10cm}|}{ Aprendizado de Máquina 1 } \\\hline
  \cellcolor{blue1}{ CH PRESENCIAL } &
    \cellcolor{blue1}{CH EAD Síncrona} &
    \cellcolor{blue1}{CH EAD Síncrona Mediada} &
    \cellcolor{blue1}{CH EAD Assíncrona} &
    \cellcolor{blue1}{CH APCC} \\\hline
  \cellcolor{blue1}{ 60 } &
    \cellcolor{blue1}{ 0 } &
    \cellcolor{blue1}{ 0 } &
    \cellcolor{blue1}{ 30 } &
    \cellcolor{blue1}{0} \\\hline
  \cellcolor{blue1}{ CH Teórica } &
    \cellcolor{blue1}{CH Prática} &
    \cellcolor{blue1}{CH Humanidades} &
    \cellcolor{blue1}{CH Extensão} &
    \cellcolor{blue1}{CH EAD TOTAL} \\\hline
  \cellcolor{blue1}{ 30 } &
    \cellcolor{blue1}{ 30 } &
    \cellcolor{blue1}{ 0 } &
    \cellcolor{blue1}{ 0 } &
    \cellcolor{blue1}{ 30 } \\\hline
  \cellcolor{blue1}{CARGA-HORÁRIA TOTAL} & \multicolumn{4}{p{10cm}|}{ \cellcolor{blue1}{ 90 } } \\\hline
  Área de conhecimento & \multicolumn{4}{p{10cm}|}{ Ciência de Dados e Inteligência Artificial } \\\hline
  Ementa & \multicolumn{4}{p{10cm}|}{ Métodos avançados de aprendizado de máquina. Modelos de ensemble, incluindo bagging, boosting e florestas aleatórias, com análise de viés e variância. Otimização e automação de modelos, abrangendo seleção de algoritmos, engenharia de características, ajuste de hiperparâmetros e AutoML. Aprendizado por reforço: conceitos básicos, agentes, ambientes e aplicações. Modelagem de séries temporais com técnicas de aprendizado de máquina. Sistemas de Recomendação. Transparência e governança da IA, incluindo viés, explicabilidade e auditoria de modelos. } \\\hline
  Competências & \multicolumn{4}{p{10cm}|}{ C4 } \\\hline


  
	\end{tabular}

\end{quadro}


\clearpage
\newpage

\begin{quadro}[ht!]
  \centering\scriptsize
\caption{Unidade Curricular Comunicação Oral e Estratégica}

\begin{tabular}{|p{2.5cm}|p{2.5cm}|p{2.5cm}|p{2.5cm}|p{2.5cm}|}\hline

  Unidade curricular & \multicolumn{4}{p{10cm}|}{ Comunicação Oral e Estratégica }\\\hline
  Código & 5.2 & Período & \multicolumn{2}{p{5cm}|}{ 5 }\\\hline
  Pré-Requisitos & \multicolumn{4}{p{10cm}|}{  } \\\hline
  \cellcolor{blue1}{ CH PRESENCIAL } &
    \cellcolor{blue1}{CH EAD Síncrona} &
    \cellcolor{blue1}{CH EAD Síncrona Mediada} &
    \cellcolor{blue1}{CH EAD Assíncrona} &
    \cellcolor{blue1}{CH APCC} \\\hline
  \cellcolor{blue1}{ 30 } &
    \cellcolor{blue1}{ 0 } &
    \cellcolor{blue1}{ 0 } &
    \cellcolor{blue1}{ 0 } &
    \cellcolor{blue1}{0} \\\hline
  \cellcolor{blue1}{ CH Teórica } &
    \cellcolor{blue1}{CH Prática} &
    \cellcolor{blue1}{CH Humanidades} &
    \cellcolor{blue1}{CH Extensão} &
    \cellcolor{blue1}{CH EAD TOTAL} \\\hline
  \cellcolor{blue1}{ 15 } &
    \cellcolor{blue1}{ 15 } &
    \cellcolor{blue1}{ 30 } &
    \cellcolor{blue1}{ 0 } &
    \cellcolor{blue1}{ 0 } \\\hline
  \cellcolor{blue1}{CARGA-HORÁRIA TOTAL} & \multicolumn{4}{p{10cm}|}{ \cellcolor{blue1}{ 30 } } \\\hline
  Área de conhecimento & \multicolumn{4}{p{10cm}|}{ Humanidades } \\\hline
  Ementa & \multicolumn{4}{p{10cm}|}{ Aspectos dos processos de produção de gêneros orais empresariais e de comunicação para apresentação em público. Planejamento, produção e avaliação de apresentação pessoal conforme as diretrizes de comunicação oral eficaz. Exposição oral de temáticas da área. Comunicação assertiva, autônoma e criativa. Aplicação de procedimentos relacionados à expressão corporal, vestuário, recursos tecnológicos e estratégias comunicativas mais adequadas ao público-alvo. } \\\hline
  Competências & \multicolumn{4}{p{10cm}|}{ C8 } \\\hline


  
	\end{tabular}

\end{quadro}


\clearpage
\newpage

\begin{quadro}[ht!]
  \centering\scriptsize
\caption{Unidade Curricular Engenharia de Software para CDIA}

\begin{tabular}{|p{2.5cm}|p{2.5cm}|p{2.5cm}|p{2.5cm}|p{2.5cm}|}\hline

  Unidade curricular & \multicolumn{4}{p{10cm}|}{ Engenharia de Software para CDIA }\\\hline
  Código & 5.3 & Período & \multicolumn{2}{p{5cm}|}{ 5 }\\\hline
  Pré-Requisitos & \multicolumn{4}{p{10cm}|}{ Programação Orientada a Objetos } \\\hline
  \cellcolor{blue1}{ CH PRESENCIAL } &
    \cellcolor{blue1}{CH EAD Síncrona} &
    \cellcolor{blue1}{CH EAD Síncrona Mediada} &
    \cellcolor{blue1}{CH EAD Assíncrona} &
    \cellcolor{blue1}{CH APCC} \\\hline
  \cellcolor{blue1}{ 30 } &
    \cellcolor{blue1}{ 0 } &
    \cellcolor{blue1}{ 0 } &
    \cellcolor{blue1}{ 30 } &
    \cellcolor{blue1}{0} \\\hline
  \cellcolor{blue1}{ CH Teórica } &
    \cellcolor{blue1}{CH Prática} &
    \cellcolor{blue1}{CH Humanidades} &
    \cellcolor{blue1}{CH Extensão} &
    \cellcolor{blue1}{CH EAD TOTAL} \\\hline
  \cellcolor{blue1}{ 30 } &
    \cellcolor{blue1}{ 30 } &
    \cellcolor{blue1}{ 0 } &
    \cellcolor{blue1}{ 0 } &
    \cellcolor{blue1}{ 30 } \\\hline
  \cellcolor{blue1}{CARGA-HORÁRIA TOTAL} & \multicolumn{4}{p{10cm}|}{ \cellcolor{blue1}{ 60 } } \\\hline
  Área de conhecimento & \multicolumn{4}{p{10cm}|}{ Computação } \\\hline
  Ementa & \multicolumn{4}{p{10cm}|}{ Fundamentos de Engenharia de Software e Requisitos: Conceitos básicos de ES, Ciclos de Vida (cascata, iterativo) e Metodologias Ágeis (Scrum, Kanban). Requisitos funcionais e não-funcionais, Requisitos específicos para Sistemas de IA (data quality, explicabilidade, latency), LGPD. Modelagem, Design e Arquitetura de Software: Diagrama de Classes, Diagrama de Componentes e Diagrama de Implantação, Princípios de Design (SOLID). Padrões de Projeto Orientados a Objetos (padrões GoF), Arquiteturas de Software (microsserviços), Componentes Essenciais de Sistemas de ML (pipelines de dados, serving).  } \\\hline
  Competências & \multicolumn{4}{p{10cm}|}{ C7 } \\\hline


  
	\end{tabular}

\end{quadro}


\clearpage
\newpage

\begin{quadro}[ht!]
  \centering\scriptsize
\caption{Unidade Curricular Metodologia de Pesquisa}

\begin{tabular}{|p{2.5cm}|p{2.5cm}|p{2.5cm}|p{2.5cm}|p{2.5cm}|}\hline

  Unidade curricular & \multicolumn{4}{p{10cm}|}{ Metodologia de Pesquisa }\\\hline
  Código & 5.4 & Período & \multicolumn{2}{p{5cm}|}{ 5 }\\\hline
  Pré-Requisitos & \multicolumn{4}{p{10cm}|}{ Período: 5 } \\\hline
  \cellcolor{blue1}{ CH PRESENCIAL } &
    \cellcolor{blue1}{CH EAD Síncrona} &
    \cellcolor{blue1}{CH EAD Síncrona Mediada} &
    \cellcolor{blue1}{CH EAD Assíncrona} &
    \cellcolor{blue1}{CH APCC} \\\hline
  \cellcolor{blue1}{ 30 } &
    \cellcolor{blue1}{ 0 } &
    \cellcolor{blue1}{ 0 } &
    \cellcolor{blue1}{ 0 } &
    \cellcolor{blue1}{0} \\\hline
  \cellcolor{blue1}{ CH Teórica } &
    \cellcolor{blue1}{CH Prática} &
    \cellcolor{blue1}{CH Humanidades} &
    \cellcolor{blue1}{CH Extensão} &
    \cellcolor{blue1}{CH EAD TOTAL} \\\hline
  \cellcolor{blue1}{ 15 } &
    \cellcolor{blue1}{ 15 } &
    \cellcolor{blue1}{ 30 } &
    \cellcolor{blue1}{ 0 } &
    \cellcolor{blue1}{ 0 } \\\hline
  \cellcolor{blue1}{CARGA-HORÁRIA TOTAL} & \multicolumn{4}{p{10cm}|}{ \cellcolor{blue1}{ 30 } } \\\hline
  Área de conhecimento & \multicolumn{4}{p{10cm}|}{ Humanidades } \\\hline
  Ementa & \multicolumn{4}{p{10cm}|}{ Fundamentos da metodologia científica: ciência e tecnologia, pesquisa e desenvolvimento tecnológico, inovação tecnológica, tipos de trabalhos acadêmicos e níveis de trabalhos de conclusão (TCC de graduação e especialização, dissertação e tese), processo de produção e comunicação científica (projeto de pesquisa, estrutura do projeto de pesquisa), estrutura de um artigo científico, comparativo entre artigo científico e trabalho de conclusão. Planejamento e desenvolvimento da pesquisa: problema de pesquisa, objetivos e hipóteses de pesquisa. Elaboração do referencial teórico: busca em bases de dados bibliográficas, técnicas de revisão da literatura (revisão bibliométrica, revisão e mapeamento sistemático), ferramentas de apoio, normas ABNT de citações e referências e escrita científica. Tipos de pesquisa: classificação das pesquisas e métodos de pesquisa, survey, estudo de caso, pesquisa-ação, pesquisa experimental (experimento e quasi-experimento). Elaboração do trabalho de conclusão de curso: revisão bibliográfica, definição de tema e objetivos (geral e específicos) e comunicação científica (refinamento do trabalho acadêmico). } \\\hline
  Competências & \multicolumn{4}{p{10cm}|}{ C7 } \\\hline


  
	\end{tabular}

\end{quadro}


\clearpage
\newpage

\begin{quadro}[ht!]
  \centering\scriptsize
\caption{Unidade Curricular Visualização de Dados}

\begin{tabular}{|p{2.5cm}|p{2.5cm}|p{2.5cm}|p{2.5cm}|p{2.5cm}|}\hline

  Unidade curricular & \multicolumn{4}{p{10cm}|}{ Visualização de Dados }\\\hline
  Código & 5.5 & Período & \multicolumn{2}{p{5cm}|}{ 5 }\\\hline
  Pré-Requisitos & \multicolumn{4}{p{10cm}|}{ Engenharia de Dados } \\\hline
  \cellcolor{blue1}{ CH PRESENCIAL } &
    \cellcolor{blue1}{CH EAD Síncrona} &
    \cellcolor{blue1}{CH EAD Síncrona Mediada} &
    \cellcolor{blue1}{CH EAD Assíncrona} &
    \cellcolor{blue1}{CH APCC} \\\hline
  \cellcolor{blue1}{ 60 } &
    \cellcolor{blue1}{ 0 } &
    \cellcolor{blue1}{ 0 } &
    \cellcolor{blue1}{ 0 } &
    \cellcolor{blue1}{0} \\\hline
  \cellcolor{blue1}{ CH Teórica } &
    \cellcolor{blue1}{CH Prática} &
    \cellcolor{blue1}{CH Humanidades} &
    \cellcolor{blue1}{CH Extensão} &
    \cellcolor{blue1}{CH EAD TOTAL} \\\hline
  \cellcolor{blue1}{ 30 } &
    \cellcolor{blue1}{ 30 } &
    \cellcolor{blue1}{ 0 } &
    \cellcolor{blue1}{ 0 } &
    \cellcolor{blue1}{ 0 } \\\hline
  \cellcolor{blue1}{CARGA-HORÁRIA TOTAL} & \multicolumn{4}{p{10cm}|}{ \cellcolor{blue1}{ 60 } } \\\hline
  Área de conhecimento & \multicolumn{4}{p{10cm}|}{ Ciência de Dados e Inteligência Artificial } \\\hline
  Ementa & \multicolumn{4}{p{10cm}|}{ Princípios de percepção visual aplicados à análise de dados; tipos de dados e mapeamento visual; seleção adequada de representações gráficas para análise exploratória. Visualização exploratória e explicativa como apoio à interpretação e comunicação de resultados. Construção de dashboards analíticos: organização da informação, layouts, interatividade e filtros. Uso de painéis dinâmicos como suporte à análise e à tomada de decisão. } \\\hline
  Competências & \multicolumn{4}{p{10cm}|}{ C5 } \\\hline


  
	\end{tabular}

\end{quadro}


\clearpage
\newpage

\begin{quadro}[ht!]
  \centering\scriptsize
\caption{Unidade Curricular Web API e Integração de Dados}

\begin{tabular}{|p{2.5cm}|p{2.5cm}|p{2.5cm}|p{2.5cm}|p{2.5cm}|}\hline

  Unidade curricular & \multicolumn{4}{p{10cm}|}{ Web API e Integração de Dados }\\\hline
  Código & 5.6 & Período & \multicolumn{2}{p{5cm}|}{ 5 }\\\hline
  Pré-Requisitos & \multicolumn{4}{p{10cm}|}{ Programação Orientada a Objetos } \\\hline
  \cellcolor{blue1}{ CH PRESENCIAL } &
    \cellcolor{blue1}{CH EAD Síncrona} &
    \cellcolor{blue1}{CH EAD Síncrona Mediada} &
    \cellcolor{blue1}{CH EAD Assíncrona} &
    \cellcolor{blue1}{CH APCC} \\\hline
  \cellcolor{blue1}{ 30 } &
    \cellcolor{blue1}{ 0 } &
    \cellcolor{blue1}{ 0 } &
    \cellcolor{blue1}{ 30 } &
    \cellcolor{blue1}{0} \\\hline
  \cellcolor{blue1}{ CH Teórica } &
    \cellcolor{blue1}{CH Prática} &
    \cellcolor{blue1}{CH Humanidades} &
    \cellcolor{blue1}{CH Extensão} &
    \cellcolor{blue1}{CH EAD TOTAL} \\\hline
  \cellcolor{blue1}{ 30 } &
    \cellcolor{blue1}{ 30 } &
    \cellcolor{blue1}{ 0 } &
    \cellcolor{blue1}{ 0 } &
    \cellcolor{blue1}{ 30 } \\\hline
  \cellcolor{blue1}{CARGA-HORÁRIA TOTAL} & \multicolumn{4}{p{10cm}|}{ \cellcolor{blue1}{ 60 } } \\\hline
  Área de conhecimento & \multicolumn{4}{p{10cm}|}{ Computação } \\\hline
  Ementa & \multicolumn{4}{p{10cm}|}{ Tecnologias de Web: HTML, HTTP, Browser, CSS, JavaScript. Programação no servidor: arquitetura cliente-servidor com HTTP, recebimento de requisições, envio de respostas, tratamento de parâmetros, cookies/sessões, conexão com banco de dados. Arquiteturas Web: arquiteturas em camadas, monolitos/microserviços, webservices/REST, filas de mensagens/fluxo de eventos e escalabilidade (estratégias de cache, redundância, replicação/fragmentação e nuvem). Segurança em aplicações web: criptografia, XSS, SQL Inject, HTTPS, chaves assimétricas e certificados digitais. } \\\hline
  Competências & \multicolumn{4}{p{10cm}|}{ C4 } \\\hline


  
	\end{tabular}

\end{quadro}


\clearpage
\newpage

\begin{quadro}[ht!]
  \centering\scriptsize
\caption{Unidade Curricular Aprendizado Profundo}

\begin{tabular}{|p{2.5cm}|p{2.5cm}|p{2.5cm}|p{2.5cm}|p{2.5cm}|}\hline

  Unidade curricular & \multicolumn{4}{p{10cm}|}{ Aprendizado Profundo }\\\hline
  Código & 6.1 & Período & \multicolumn{2}{p{5cm}|}{ 6 }\\\hline
  Pré-Requisitos & \multicolumn{4}{p{10cm}|}{ Aprendizado de Máquina 1 } \\\hline
  \cellcolor{blue1}{ CH PRESENCIAL } &
    \cellcolor{blue1}{CH EAD Síncrona} &
    \cellcolor{blue1}{CH EAD Síncrona Mediada} &
    \cellcolor{blue1}{CH EAD Assíncrona} &
    \cellcolor{blue1}{CH APCC} \\\hline
  \cellcolor{blue1}{ 60 } &
    \cellcolor{blue1}{ 0 } &
    \cellcolor{blue1}{ 0 } &
    \cellcolor{blue1}{ 0 } &
    \cellcolor{blue1}{0} \\\hline
  \cellcolor{blue1}{ CH Teórica } &
    \cellcolor{blue1}{CH Prática} &
    \cellcolor{blue1}{CH Humanidades} &
    \cellcolor{blue1}{CH Extensão} &
    \cellcolor{blue1}{CH EAD TOTAL} \\\hline
  \cellcolor{blue1}{ 30 } &
    \cellcolor{blue1}{ 30 } &
    \cellcolor{blue1}{ 0 } &
    \cellcolor{blue1}{ 0 } &
    \cellcolor{blue1}{ 0 } \\\hline
  \cellcolor{blue1}{CARGA-HORÁRIA TOTAL} & \multicolumn{4}{p{10cm}|}{ \cellcolor{blue1}{ 60 } } \\\hline
  Área de conhecimento & \multicolumn{4}{p{10cm}|}{ Ciência de Dados e Inteligência Artificial } \\\hline
  Ementa & \multicolumn{4}{p{10cm}|}{ Fundamentos de Redes Neurais Artificiais: Aprendizagem supervisionada e não supervisionada; neurônio artificial, funções de ativação, função objetivo, o algoritmo de Descida de Gradiente e Regularizações. Aprendizagem profunda (Deep Learning): Redes Neurais Recorrentes; Redes Neurais Convolucionais; aplicações de Aprendizagem Profunda. } \\\hline
  Competências & \multicolumn{4}{p{10cm}|}{ C4 } \\\hline


  
	\end{tabular}

\end{quadro}


\clearpage
\newpage

\begin{quadro}[ht!]
  \centering\scriptsize
\caption{Unidade Curricular Big Data e Computação em Nuvem}

\begin{tabular}{|p{2.5cm}|p{2.5cm}|p{2.5cm}|p{2.5cm}|p{2.5cm}|}\hline

  Unidade curricular & \multicolumn{4}{p{10cm}|}{ Big Data e Computação em Nuvem }\\\hline
  Código & 6.2 & Período & \multicolumn{2}{p{5cm}|}{ 6 }\\\hline
  Pré-Requisitos & \multicolumn{4}{p{10cm}|}{ Engenharia de Dados } \\\hline
  \cellcolor{blue1}{ CH PRESENCIAL } &
    \cellcolor{blue1}{CH EAD Síncrona} &
    \cellcolor{blue1}{CH EAD Síncrona Mediada} &
    \cellcolor{blue1}{CH EAD Assíncrona} &
    \cellcolor{blue1}{CH APCC} \\\hline
  \cellcolor{blue1}{ 30 } &
    \cellcolor{blue1}{ 0 } &
    \cellcolor{blue1}{ 0 } &
    \cellcolor{blue1}{ 30 } &
    \cellcolor{blue1}{0} \\\hline
  \cellcolor{blue1}{ CH Teórica } &
    \cellcolor{blue1}{CH Prática} &
    \cellcolor{blue1}{CH Humanidades} &
    \cellcolor{blue1}{CH Extensão} &
    \cellcolor{blue1}{CH EAD TOTAL} \\\hline
  \cellcolor{blue1}{ 30 } &
    \cellcolor{blue1}{ 30 } &
    \cellcolor{blue1}{ 0 } &
    \cellcolor{blue1}{ 0 } &
    \cellcolor{blue1}{ 30 } \\\hline
  \cellcolor{blue1}{CARGA-HORÁRIA TOTAL} & \multicolumn{4}{p{10cm}|}{ \cellcolor{blue1}{ 60 } } \\\hline
  Área de conhecimento & \multicolumn{4}{p{10cm}|}{ Ciência de Dados e Inteligência Artificial } \\\hline
  Ementa & \multicolumn{4}{p{10cm}|}{ Fundamentos e Infraestrutura de Big Data: Definição, evolução e os 5 Vs (Volume, Velocidade, Variedade, Veracidade e Valor), Sistemas Distribuídos e Paralelos para BigData, Ecossistema Hadoop, Apache Spark, Arquitetura Data Lake (conceitos, comparação com Data Warehouse). Computação em Nuvem: Conceitos de Computação em Nuvem, Serviços de armazenamento e processamento distribuído na nuvem. BigData, Nuvem e Aplicações de IA: Plataformas de Machine Learning (MLOps) e serviços de IA gerenciados, aplicações de Big Data, Computação em Nuvem e Inteligência Artificial. } \\\hline
  Competências & \multicolumn{4}{p{10cm}|}{ C3 } \\\hline


  
	\end{tabular}

\end{quadro}


\clearpage
\newpage

\begin{quadro}[ht!]
  \centering\scriptsize
\caption{Unidade Curricular Desenvolvimento de Novos Negócios}

\begin{tabular}{|p{2.5cm}|p{2.5cm}|p{2.5cm}|p{2.5cm}|p{2.5cm}|}\hline

  Unidade curricular & \multicolumn{4}{p{10cm}|}{ Desenvolvimento de Novos Negócios }\\\hline
  Código & 6.3 & Período & \multicolumn{2}{p{5cm}|}{ 6 }\\\hline
  Pré-Requisitos & \multicolumn{4}{p{10cm}|}{  } \\\hline
  \cellcolor{blue1}{ CH PRESENCIAL } &
    \cellcolor{blue1}{CH EAD Síncrona} &
    \cellcolor{blue1}{CH EAD Síncrona Mediada} &
    \cellcolor{blue1}{CH EAD Assíncrona} &
    \cellcolor{blue1}{CH APCC} \\\hline
  \cellcolor{blue1}{ 0 } &
    \cellcolor{blue1}{ 0 } &
    \cellcolor{blue1}{ 30 } &
    \cellcolor{blue1}{ 0 } &
    \cellcolor{blue1}{0} \\\hline
  \cellcolor{blue1}{ CH Teórica } &
    \cellcolor{blue1}{CH Prática} &
    \cellcolor{blue1}{CH Humanidades} &
    \cellcolor{blue1}{CH Extensão} &
    \cellcolor{blue1}{CH EAD TOTAL} \\\hline
  \cellcolor{blue1}{ 15 } &
    \cellcolor{blue1}{ 15 } &
    \cellcolor{blue1}{ 30 } &
    \cellcolor{blue1}{ 0 } &
    \cellcolor{blue1}{ 30 } \\\hline
  \cellcolor{blue1}{CARGA-HORÁRIA TOTAL} & \multicolumn{4}{p{10cm}|}{ \cellcolor{blue1}{ 30 } } \\\hline
  Área de conhecimento & \multicolumn{4}{p{10cm}|}{ Humanidades } \\\hline
  Ementa & \multicolumn{4}{p{10cm}|}{ Características e Comportamento Empreendedor. Intraempreendedorismo. Identificação de Oportunidades de negócios. Elaboração de um Plano de Negócios. Fontes de investimento para novos negócios.  } \\\hline
  Competências & \multicolumn{4}{p{10cm}|}{ C7 } \\\hline


  
	\end{tabular}

\end{quadro}


\clearpage
\newpage

\begin{quadro}[ht!]
  \centering\scriptsize
\caption{Unidade Curricular Implantação e Operação de Modelos}

\begin{tabular}{|p{2.5cm}|p{2.5cm}|p{2.5cm}|p{2.5cm}|p{2.5cm}|}\hline

  Unidade curricular & \multicolumn{4}{p{10cm}|}{ Implantação e Operação de Modelos }\\\hline
  Código & 6.4 & Período & \multicolumn{2}{p{5cm}|}{ 6 }\\\hline
  Pré-Requisitos & \multicolumn{4}{p{10cm}|}{ Engenharia de Software para CDIA } \\\hline
  \cellcolor{blue1}{ CH PRESENCIAL } &
    \cellcolor{blue1}{CH EAD Síncrona} &
    \cellcolor{blue1}{CH EAD Síncrona Mediada} &
    \cellcolor{blue1}{CH EAD Assíncrona} &
    \cellcolor{blue1}{CH APCC} \\\hline
  \cellcolor{blue1}{ 60 } &
    \cellcolor{blue1}{ 0 } &
    \cellcolor{blue1}{ 0 } &
    \cellcolor{blue1}{ 0 } &
    \cellcolor{blue1}{0} \\\hline
  \cellcolor{blue1}{ CH Teórica } &
    \cellcolor{blue1}{CH Prática} &
    \cellcolor{blue1}{CH Humanidades} &
    \cellcolor{blue1}{CH Extensão} &
    \cellcolor{blue1}{CH EAD TOTAL} \\\hline
  \cellcolor{blue1}{ 30 } &
    \cellcolor{blue1}{ 30 } &
    \cellcolor{blue1}{ 0 } &
    \cellcolor{blue1}{ 0 } &
    \cellcolor{blue1}{ 0 } \\\hline
  \cellcolor{blue1}{CARGA-HORÁRIA TOTAL} & \multicolumn{4}{p{10cm}|}{ \cellcolor{blue1}{ 60 } } \\\hline
  Área de conhecimento & \multicolumn{4}{p{10cm}|}{ Ciência de Dados e Inteligência Artificial } \\\hline
  Ementa & \multicolumn{4}{p{10cm}|}{ Gerência de Configuração: ciclo de vida da GC. Conceitos de repositórios, referenciais (baselines), e workspace. Versionamento de Código (Git e Git Flow). GC aplicada a ML: Versionamento de Dados e Modelos (DVC). Teste, Qualidade e MLOps: Qualidade de Software, Conceitos de Verificação/Validação e técnicas de Inspeção de código, Estratégias e Tipos de Testes (Unidade, Integração, Sistema, Aceitação), Testes específicos para ML (validação de dados, modelo e bias), Integração Contínua (CI) e Entrega Contínua (CD), Princípios de MLOps, Monitoramento de modelos (data drift, model drift). } \\\hline
  Competências & \multicolumn{4}{p{10cm}|}{ C5 } \\\hline


  
	\end{tabular}

\end{quadro}


\clearpage
\newpage

\begin{quadro}[ht!]
  \centering\scriptsize
\caption{Unidade Curricular Oficina de Integração 1}

\begin{tabular}{|p{2.5cm}|p{2.5cm}|p{2.5cm}|p{2.5cm}|p{2.5cm}|}\hline

  Unidade curricular & \multicolumn{4}{p{10cm}|}{ Oficina de Integração 1 }\\\hline
  Código & 6.5 & Período & \multicolumn{2}{p{5cm}|}{ 6 }\\\hline
  Pré-Requisitos & \multicolumn{4}{p{10cm}|}{ Período: 6 } \\\hline
  \cellcolor{blue1}{ CH PRESENCIAL } &
    \cellcolor{blue1}{CH EAD Síncrona} &
    \cellcolor{blue1}{CH EAD Síncrona Mediada} &
    \cellcolor{blue1}{CH EAD Assíncrona} &
    \cellcolor{blue1}{CH APCC} \\\hline
  \cellcolor{blue1}{ 120 } &
    \cellcolor{blue1}{ 0 } &
    \cellcolor{blue1}{ 0 } &
    \cellcolor{blue1}{ 0 } &
    \cellcolor{blue1}{0} \\\hline
  \cellcolor{blue1}{ CH Teórica } &
    \cellcolor{blue1}{CH Prática} &
    \cellcolor{blue1}{CH Humanidades} &
    \cellcolor{blue1}{CH Extensão} &
    \cellcolor{blue1}{CH EAD TOTAL} \\\hline
  \cellcolor{blue1}{ 0 } &
    \cellcolor{blue1}{ 0 } &
    \cellcolor{blue1}{ 0 } &
    \cellcolor{blue1}{ 120 } &
    \cellcolor{blue1}{ 0 } \\\hline
  \cellcolor{blue1}{CARGA-HORÁRIA TOTAL} & \multicolumn{4}{p{10cm}|}{ \cellcolor{blue1}{ 120 } } \\\hline
  Área de conhecimento & \multicolumn{4}{p{10cm}|}{ Extensão } \\\hline
  Ementa & \multicolumn{4}{p{10cm}|}{ Integração dos conhecimentos das disciplinas da área de Computação do primeiro, segundo e terceiro, quarto e quinto períodos. Aplicação desses conhecimentos em todas as etapas do desenvolvimento de um sistema computacional que resolva um problema real da comunidade externa a UTFPR, nas temáticas desenvolvidas no campus das áreas de Sustentabilidade e IA para o bem. } \\\hline
  Competências & \multicolumn{4}{p{10cm}|}{ C6 } \\\hline


  
	\end{tabular}

\end{quadro}


\clearpage
\newpage

\begin{quadro}[ht!]
  \centering\scriptsize
\caption{Unidade Curricular Processamento de Linguagem Natural}

\begin{tabular}{|p{2.5cm}|p{2.5cm}|p{2.5cm}|p{2.5cm}|p{2.5cm}|}\hline

  Unidade curricular & \multicolumn{4}{p{10cm}|}{ Processamento de Linguagem Natural }\\\hline
  Código & 6.6 & Período & \multicolumn{2}{p{5cm}|}{ 6 }\\\hline
  Pré-Requisitos & \multicolumn{4}{p{10cm}|}{ Aprendizado de Máquina 1 } \\\hline
  \cellcolor{blue1}{ CH PRESENCIAL } &
    \cellcolor{blue1}{CH EAD Síncrona} &
    \cellcolor{blue1}{CH EAD Síncrona Mediada} &
    \cellcolor{blue1}{CH EAD Assíncrona} &
    \cellcolor{blue1}{CH APCC} \\\hline
  \cellcolor{blue1}{ 60 } &
    \cellcolor{blue1}{ 0 } &
    \cellcolor{blue1}{ 0 } &
    \cellcolor{blue1}{ 0 } &
    \cellcolor{blue1}{0} \\\hline
  \cellcolor{blue1}{ CH Teórica } &
    \cellcolor{blue1}{CH Prática} &
    \cellcolor{blue1}{CH Humanidades} &
    \cellcolor{blue1}{CH Extensão} &
    \cellcolor{blue1}{CH EAD TOTAL} \\\hline
  \cellcolor{blue1}{ 30 } &
    \cellcolor{blue1}{ 30 } &
    \cellcolor{blue1}{ 0 } &
    \cellcolor{blue1}{ 0 } &
    \cellcolor{blue1}{ 0 } \\\hline
  \cellcolor{blue1}{CARGA-HORÁRIA TOTAL} & \multicolumn{4}{p{10cm}|}{ \cellcolor{blue1}{ 60 } } \\\hline
  Área de conhecimento & \multicolumn{4}{p{10cm}|}{ Ciência de Dados e Inteligência Artificial } \\\hline
  Ementa & \multicolumn{4}{p{10cm}|}{ Fundamentos e aplicações de Processamento de Linguagem Natural. Pré-processamento de textos (córpus), incluindo tokenização, normalização, lematização e reconhecimento de entidades. Representação textual baseada em frequência, n-gramas, embeddings e extração de características. Aplicações clássicas de PLN em recuperação de informação, classificação, sumarização, clusterização e análise de tópicos. } \\\hline
  Competências & \multicolumn{4}{p{10cm}|}{ C4 } \\\hline


  
	\end{tabular}

\end{quadro}


\clearpage
\newpage

\begin{quadro}[ht!]
  \centering\scriptsize
\caption{Unidade Curricular Computação de Alto Desempenho}

\begin{tabular}{|p{2.5cm}|p{2.5cm}|p{2.5cm}|p{2.5cm}|p{2.5cm}|}\hline

  Unidade curricular & \multicolumn{4}{p{10cm}|}{ Computação de Alto Desempenho }\\\hline
  Código & 7.1 & Período & \multicolumn{2}{p{5cm}|}{ 7 }\\\hline
  Pré-Requisitos & \multicolumn{4}{p{10cm}|}{ Sistemas de Computação e Comunicação de Dados } \\\hline
  \cellcolor{blue1}{ CH PRESENCIAL } &
    \cellcolor{blue1}{CH EAD Síncrona} &
    \cellcolor{blue1}{CH EAD Síncrona Mediada} &
    \cellcolor{blue1}{CH EAD Assíncrona} &
    \cellcolor{blue1}{CH APCC} \\\hline
  \cellcolor{blue1}{ 60 } &
    \cellcolor{blue1}{ 0 } &
    \cellcolor{blue1}{ 0 } &
    \cellcolor{blue1}{ 0 } &
    \cellcolor{blue1}{0} \\\hline
  \cellcolor{blue1}{ CH Teórica } &
    \cellcolor{blue1}{CH Prática} &
    \cellcolor{blue1}{CH Humanidades} &
    \cellcolor{blue1}{CH Extensão} &
    \cellcolor{blue1}{CH EAD TOTAL} \\\hline
  \cellcolor{blue1}{ 30 } &
    \cellcolor{blue1}{ 30 } &
    \cellcolor{blue1}{ 0 } &
    \cellcolor{blue1}{ 0 } &
    \cellcolor{blue1}{ 0 } \\\hline
  \cellcolor{blue1}{CARGA-HORÁRIA TOTAL} & \multicolumn{4}{p{10cm}|}{ \cellcolor{blue1}{ 60 } } \\\hline
  Área de conhecimento & \multicolumn{4}{p{10cm}|}{ Ciência de Dados e Inteligência Artificial } \\\hline
  Ementa & \multicolumn{4}{p{10cm}|}{ Fundamentos Teóricos e Avaliação de Desempenho: Introdução ao paralelismo (motivação, taxonomia), arquiteturas de hardware (memória compartilhada vs. distribuída), e os princípios de projeto de algoritmos. Programação para Memória Compartilhada e Distribuída: Técnicas de programação para ambientes multicore (uma única máquina com múltiplos núcleos) usando OpenMP, com foco em threads e sincronização. Aceleração de Computação (GPUs e Arquiteturas Heterogêneas): Arquitetura de GPUs e o modelo de programação em massa (como CUDA ou OpenCL). } \\\hline
  Competências & \multicolumn{4}{p{10cm}|}{ C3 } \\\hline


  
	\end{tabular}

\end{quadro}


\clearpage
\newpage

\begin{quadro}[ht!]
  \centering\scriptsize
\caption{Unidade Curricular Inteligência Artificial Generativa}

\begin{tabular}{|p{2.5cm}|p{2.5cm}|p{2.5cm}|p{2.5cm}|p{2.5cm}|}\hline

  Unidade curricular & \multicolumn{4}{p{10cm}|}{ Inteligência Artificial Generativa }\\\hline
  Código & 7.2 & Período & \multicolumn{2}{p{5cm}|}{ 7 }\\\hline
  Pré-Requisitos & \multicolumn{4}{p{10cm}|}{ Aprendizado Profundo } \\\hline
  \cellcolor{blue1}{ CH PRESENCIAL } &
    \cellcolor{blue1}{CH EAD Síncrona} &
    \cellcolor{blue1}{CH EAD Síncrona Mediada} &
    \cellcolor{blue1}{CH EAD Assíncrona} &
    \cellcolor{blue1}{CH APCC} \\\hline
  \cellcolor{blue1}{ 0 } &
    \cellcolor{blue1}{ 0 } &
    \cellcolor{blue1}{ 30 } &
    \cellcolor{blue1}{ 30 } &
    \cellcolor{blue1}{0} \\\hline
  \cellcolor{blue1}{ CH Teórica } &
    \cellcolor{blue1}{CH Prática} &
    \cellcolor{blue1}{CH Humanidades} &
    \cellcolor{blue1}{CH Extensão} &
    \cellcolor{blue1}{CH EAD TOTAL} \\\hline
  \cellcolor{blue1}{ 30 } &
    \cellcolor{blue1}{ 30 } &
    \cellcolor{blue1}{ 0 } &
    \cellcolor{blue1}{ 0 } &
    \cellcolor{blue1}{ 60 } \\\hline
  \cellcolor{blue1}{CARGA-HORÁRIA TOTAL} & \multicolumn{4}{p{10cm}|}{ \cellcolor{blue1}{ 60 } } \\\hline
  Área de conhecimento & \multicolumn{4}{p{10cm}|}{ Ciência de Dados e Inteligência Artificial } \\\hline
  Ementa & \multicolumn{4}{p{10cm}|}{ Fundamentos de IA Generativa aplicada à linguagem natural. Arquiteturas modernas baseadas em modelos do tipo Transformer e Large Language Models. Transferência de aprendizado e uso de modelos pré-treinados. Técnicas de interação e controle de modelos generativos, incluindo formulação de prompts e estratégias de orientação de respostas. Aplicações em compreensão, geração de texto, recuperação de informação, sumarização e tradução, com discussão de limitações, riscos e aspectos éticos. Embeddings, Busca Vetorial, Banco de Dados Vetoriais, RAG. Multimodal Embeddings. } \\\hline
  Competências & \multicolumn{4}{p{10cm}|}{ C4 } \\\hline


  
	\end{tabular}

\end{quadro}


\clearpage
\newpage

\begin{quadro}[ht!]
  \centering\scriptsize
\caption{Unidade Curricular Oficina de Integração 2}

\begin{tabular}{|p{2.5cm}|p{2.5cm}|p{2.5cm}|p{2.5cm}|p{2.5cm}|}\hline

  Unidade curricular & \multicolumn{4}{p{10cm}|}{ Oficina de Integração 2 }\\\hline
  Código & 7.3 & Período & \multicolumn{2}{p{5cm}|}{ 7 }\\\hline
  Pré-Requisitos & \multicolumn{4}{p{10cm}|}{ Período: 7 } \\\hline
  \cellcolor{blue1}{ CH PRESENCIAL } &
    \cellcolor{blue1}{CH EAD Síncrona} &
    \cellcolor{blue1}{CH EAD Síncrona Mediada} &
    \cellcolor{blue1}{CH EAD Assíncrona} &
    \cellcolor{blue1}{CH APCC} \\\hline
  \cellcolor{blue1}{ 120 } &
    \cellcolor{blue1}{ 0 } &
    \cellcolor{blue1}{ 0 } &
    \cellcolor{blue1}{ 0 } &
    \cellcolor{blue1}{0} \\\hline
  \cellcolor{blue1}{ CH Teórica } &
    \cellcolor{blue1}{CH Prática} &
    \cellcolor{blue1}{CH Humanidades} &
    \cellcolor{blue1}{CH Extensão} &
    \cellcolor{blue1}{CH EAD TOTAL} \\\hline
  \cellcolor{blue1}{ 0 } &
    \cellcolor{blue1}{ 0 } &
    \cellcolor{blue1}{ 0 } &
    \cellcolor{blue1}{ 120 } &
    \cellcolor{blue1}{ 0 } \\\hline
  \cellcolor{blue1}{CARGA-HORÁRIA TOTAL} & \multicolumn{4}{p{10cm}|}{ \cellcolor{blue1}{ 120 } } \\\hline
  Área de conhecimento & \multicolumn{4}{p{10cm}|}{ Extensão } \\\hline
  Ementa & \multicolumn{4}{p{10cm}|}{ Integração dos conhecimentos das disciplinas da área de Ciência de Dados e Inteligência Artificial do primeiro, segundo e terceiro, quarto, quinto e sexto períodos. Aplicação desses conhecimentos em todas as etapas do desenvolvimento de um sistema computacional que resolva um problema real da comunidade externa a UTFPR, nas temáticas desenvolvidas no campus das áreas de Sustentabilidade e IA para o bem. } \\\hline
  Competências & \multicolumn{4}{p{10cm}|}{ C6 } \\\hline


  
	\end{tabular}

\end{quadro}


\clearpage
\newpage

\begin{quadro}[ht!]
  \centering\scriptsize
\caption{Unidade Curricular Sistemas Multiagentes}

\begin{tabular}{|p{2.5cm}|p{2.5cm}|p{2.5cm}|p{2.5cm}|p{2.5cm}|}\hline

  Unidade curricular & \multicolumn{4}{p{10cm}|}{ Sistemas Multiagentes }\\\hline
  Código & 7.4 & Período & \multicolumn{2}{p{5cm}|}{ 7 }\\\hline
  Pré-Requisitos & \multicolumn{4}{p{10cm}|}{ Fundamentos de Inteligência Artificial } \\\hline
  \cellcolor{blue1}{ CH PRESENCIAL } &
    \cellcolor{blue1}{CH EAD Síncrona} &
    \cellcolor{blue1}{CH EAD Síncrona Mediada} &
    \cellcolor{blue1}{CH EAD Assíncrona} &
    \cellcolor{blue1}{CH APCC} \\\hline
  \cellcolor{blue1}{ 60 } &
    \cellcolor{blue1}{ 0 } &
    \cellcolor{blue1}{ 0 } &
    \cellcolor{blue1}{ 0 } &
    \cellcolor{blue1}{0} \\\hline
  \cellcolor{blue1}{ CH Teórica } &
    \cellcolor{blue1}{CH Prática} &
    \cellcolor{blue1}{CH Humanidades} &
    \cellcolor{blue1}{CH Extensão} &
    \cellcolor{blue1}{CH EAD TOTAL} \\\hline
  \cellcolor{blue1}{ 30 } &
    \cellcolor{blue1}{ 30 } &
    \cellcolor{blue1}{ 0 } &
    \cellcolor{blue1}{ 0 } &
    \cellcolor{blue1}{ 0 } \\\hline
  \cellcolor{blue1}{CARGA-HORÁRIA TOTAL} & \multicolumn{4}{p{10cm}|}{ \cellcolor{blue1}{ 60 } } \\\hline
  Área de conhecimento & \multicolumn{4}{p{10cm}|}{ Ciência de Dados e Inteligência Artificial } \\\hline
  Ementa & \multicolumn{4}{p{10cm}|}{ Definições, propriedades e classificações de agentes inteligentes (tipologias de agentes, arquitetura BDI); Arquiteturas Cognitivas e Deliberativas (arquiteturas simbólicas e híbridas, agentes de resolução de problemas); Sistemas Multiagentes com Agentes Colaborativos, Agentes de Recuperação de Informação, Agentes Credíveis e Caracteres Sintéticos; Engenharia de Software Orientada a Agentes (ESOA); Agentes baseados em Aprendizado; Modelagem e Simulação. } \\\hline
  Competências & \multicolumn{4}{p{10cm}|}{ C5 } \\\hline


  
	\end{tabular}

\end{quadro}


\clearpage
\newpage

\begin{quadro}[ht!]
  \centering\scriptsize
\caption{Unidade Curricular TCC 1}

\begin{tabular}{|p{2.5cm}|p{2.5cm}|p{2.5cm}|p{2.5cm}|p{2.5cm}|}\hline

  Unidade curricular & \multicolumn{4}{p{10cm}|}{ TCC 1 }\\\hline
  Código & 7.5 & Período & \multicolumn{2}{p{5cm}|}{ 7 }\\\hline
  Pré-Requisitos & \multicolumn{4}{p{10cm}|}{ Metodologia de Pesquisa } \\\hline
  \cellcolor{blue1}{ CH PRESENCIAL } &
    \cellcolor{blue1}{CH EAD Síncrona} &
    \cellcolor{blue1}{CH EAD Síncrona Mediada} &
    \cellcolor{blue1}{CH EAD Assíncrona} &
    \cellcolor{blue1}{CH APCC} \\\hline
  \cellcolor{blue1}{ 0 } &
    \cellcolor{blue1}{ 0 } &
    \cellcolor{blue1}{ 30 } &
    \cellcolor{blue1}{ 0 } &
    \cellcolor{blue1}{0} \\\hline
  \cellcolor{blue1}{ CH Teórica } &
    \cellcolor{blue1}{CH Prática} &
    \cellcolor{blue1}{CH Humanidades} &
    \cellcolor{blue1}{CH Extensão} &
    \cellcolor{blue1}{CH EAD TOTAL} \\\hline
  \cellcolor{blue1}{ 15 } &
    \cellcolor{blue1}{ 15 } &
    \cellcolor{blue1}{ 0 } &
    \cellcolor{blue1}{ 0 } &
    \cellcolor{blue1}{ 30 } \\\hline
  \cellcolor{blue1}{CARGA-HORÁRIA TOTAL} & \multicolumn{4}{p{10cm}|}{ \cellcolor{blue1}{ 30 } } \\\hline
  Área de conhecimento & \multicolumn{4}{p{10cm}|}{ Trabalho de Conclusão de Curso } \\\hline
  Ementa & \multicolumn{4}{p{10cm}|}{ Integralizar as competências adquiridas durante o curso, aplicando-as no desenvolvimento da proposta do trabalho de conclusão, sob a orientação de um professor. } \\\hline
  Competências & \multicolumn{4}{p{10cm}|}{  } \\\hline


  
	\end{tabular}

\end{quadro}


\clearpage
\newpage

\begin{quadro}[ht!]
  \centering\scriptsize
\caption{Unidade Curricular Visão Computacional}

\begin{tabular}{|p{2.5cm}|p{2.5cm}|p{2.5cm}|p{2.5cm}|p{2.5cm}|}\hline

  Unidade curricular & \multicolumn{4}{p{10cm}|}{ Visão Computacional }\\\hline
  Código & 7.6 & Período & \multicolumn{2}{p{5cm}|}{ 7 }\\\hline
  Pré-Requisitos & \multicolumn{4}{p{10cm}|}{ Aprendizado Profundo } \\\hline
  \cellcolor{blue1}{ CH PRESENCIAL } &
    \cellcolor{blue1}{CH EAD Síncrona} &
    \cellcolor{blue1}{CH EAD Síncrona Mediada} &
    \cellcolor{blue1}{CH EAD Assíncrona} &
    \cellcolor{blue1}{CH APCC} \\\hline
  \cellcolor{blue1}{ 60 } &
    \cellcolor{blue1}{ 0 } &
    \cellcolor{blue1}{ 0 } &
    \cellcolor{blue1}{ 0 } &
    \cellcolor{blue1}{0} \\\hline
  \cellcolor{blue1}{ CH Teórica } &
    \cellcolor{blue1}{CH Prática} &
    \cellcolor{blue1}{CH Humanidades} &
    \cellcolor{blue1}{CH Extensão} &
    \cellcolor{blue1}{CH EAD TOTAL} \\\hline
  \cellcolor{blue1}{ 30 } &
    \cellcolor{blue1}{ 30 } &
    \cellcolor{blue1}{ 0 } &
    \cellcolor{blue1}{ 0 } &
    \cellcolor{blue1}{ 0 } \\\hline
  \cellcolor{blue1}{CARGA-HORÁRIA TOTAL} & \multicolumn{4}{p{10cm}|}{ \cellcolor{blue1}{ 60 } } \\\hline
  Área de conhecimento & \multicolumn{4}{p{10cm}|}{ Ciência de Dados e Inteligência Artificial } \\\hline
  Ementa & \multicolumn{4}{p{10cm}|}{ Fundamentos de processamento de imagens: representação de imagens digitais, sistema visual humano, composição da cor, digitalização, modelo de imagens, métodos de espaço de estados, resolução espacial, profundidade, amostragem e quantização. Transformações de intensidade e filtragem espacial : processamento de histograma, transformação da escala de cinza, convolução, filtros espaciais de suavização e aguçamento. Realce de imagens baseado em cores e transformações morfológicas: sistemas de cor, transformação pseudocor, realce com transformação HSV, dilação, erosão, abertura e fechamento. Segmentação e descrição: limiarização local e global, crescimento de regiões, detecção de bordas, descritores de cor, forma e textura. } \\\hline
  Competências & \multicolumn{4}{p{10cm}|}{ C4 } \\\hline


  
	\end{tabular}

\end{quadro}


\clearpage
\newpage

\begin{quadro}[ht!]
  \centering\scriptsize
\caption{Unidade Curricular Ética e Legislação para CDIA}

\begin{tabular}{|p{2.5cm}|p{2.5cm}|p{2.5cm}|p{2.5cm}|p{2.5cm}|}\hline

  Unidade curricular & \multicolumn{4}{p{10cm}|}{ Ética e Legislação para CDIA }\\\hline
  Código & 8.1 & Período & \multicolumn{2}{p{5cm}|}{ 8 }\\\hline
  Pré-Requisitos & \multicolumn{4}{p{10cm}|}{  } \\\hline
  \cellcolor{blue1}{ CH PRESENCIAL } &
    \cellcolor{blue1}{CH EAD Síncrona} &
    \cellcolor{blue1}{CH EAD Síncrona Mediada} &
    \cellcolor{blue1}{CH EAD Assíncrona} &
    \cellcolor{blue1}{CH APCC} \\\hline
  \cellcolor{blue1}{ 0 } &
    \cellcolor{blue1}{ 0 } &
    \cellcolor{blue1}{ 30 } &
    \cellcolor{blue1}{ 0 } &
    \cellcolor{blue1}{0} \\\hline
  \cellcolor{blue1}{ CH Teórica } &
    \cellcolor{blue1}{CH Prática} &
    \cellcolor{blue1}{CH Humanidades} &
    \cellcolor{blue1}{CH Extensão} &
    \cellcolor{blue1}{CH EAD TOTAL} \\\hline
  \cellcolor{blue1}{ 30 } &
    \cellcolor{blue1}{ 0 } &
    \cellcolor{blue1}{ 30 } &
    \cellcolor{blue1}{ 0 } &
    \cellcolor{blue1}{ 30 } \\\hline
  \cellcolor{blue1}{CARGA-HORÁRIA TOTAL} & \multicolumn{4}{p{10cm}|}{ \cellcolor{blue1}{ 30 } } \\\hline
  Área de conhecimento & \multicolumn{4}{p{10cm}|}{ Humanidades } \\\hline
  Ementa & \multicolumn{4}{p{10cm}|}{ Fundamentos éticos aplicados à tecnologia, princípios de ética em ciência de dados e inteligência artificial, responsabilidade social, transparência, explicabilidade e prestação de contas em sistemas algorítmicos. Conceitos de privacidade, dados pessoais e sensíveis, governança de dados, segurança da informação e boas práticas no ciclo de vida dos dados. Lei Geral de Proteção de Dados (LGPD), fundamentos e princípios, direitos dos titulares, deveres dos agentes de tratamento, bases legais, sanções e impactos no desenvolvimento de sistemas baseados em dados e IA; noções de regulamentações nacionais e internacionais relacionadas à IA. Discriminação, justiça e equidade, uso responsável de modelos preditivos, análise de casos reais, avaliação de riscos éticos e legais, e diretrizes para o desenvolvimento, implantação e uso responsável de soluções em ciência de dados e inteligência artificial. } \\\hline
  Competências & \multicolumn{4}{p{10cm}|}{ C6 } \\\hline


  
	\end{tabular}

\end{quadro}


\clearpage
\newpage

